\documentclass{situationist-preprint}
\usepackage{longtable}
\usepackage{array}
\usepackage{needspace}

\title{Appendix: Empirical Studies of Boredom}
\msid{EXI-10-b}
\msfamily{EXI-10}
\msdiscipline{Meaning, Death and the Absurd}
\mstype{Companion manuscript}
\mskeywords{boredom; boredom proneness; state boredom; flow; meaning; function of emotions; annotated bibliography}

\newcommand{\stn}{58}
\newcommand{\stnemp}{44}
\newcommand{\stntheo}{14}
\newcommand{\staS}{19}
\newcommand{\staSandC}{2}
\newcommand{\staN}{37}
\newcommand{\staC}{0}
\newcommand{\stbS}{24}
\newcommand{\stbSandC}{3}
\newcommand{\stbN}{31}
\newcommand{\stbC}{0}
\newcommand{\stbsS}{23}
\newcommand{\stbsSandC}{4}
\newcommand{\stbsN}{26}
\newcommand{\stbsC}{5}
\newcommand{\stymin}{1954}
\newcommand{\stymax}{2021}
\newcommand{\stoverpos}{1}
\newcommand{\stoverposemp}{0}
\newcommand{\stoverneg}{2}


\begin{document}
\maketitle

\begin{abstract}
EXI-10-a argues that boredom signals a mismatch between one's capacities and one's activities, and that it has the adaptive function of moving the agent toward activities that match them. Both claims answer to empirical work. We survey that work and assess its bearing. The written part sets out how the principal measures, definitions and experimental paradigms bear on the mismatch thesis and on the function thesis, and where they fall short. The appendix is an annotated bibliography of \stn{} studies and theoretical works from \stymin{} to \stymax{}, each with its design, sample, measure and main finding, and with its bearing on each thesis computed from explicit codes by fixed rules. The cross-tabulations show that the evidence for the function thesis is mostly evidence about motivation, that the meaning and under-challenge clauses of the mismatch thesis are well supported, and that the over-challenge clause has the weakest support and one standing challenge, the flow model's claim that over-challenge produces anxiety rather than boredom.
\end{abstract}

\section{Introduction}

EXI-10 holds that boredom signals a mismatch between one's capacities and one's activities, and that it has an adaptive function. EXI-10-a defends both claims philosophically: boredom is an affective state with content (its Definition~2.1 and Lemma~3.1), the mismatch it represents has three forms, namely under-challenge, over-challenge without traction, and meaninglessness (its Definition~2.2), and the motivation it produces tends to restore the match, so that, given that engagement pays (its Assumption~2.4), the disposition to boredom has the selected-effect function of moving the agent toward matching activities (its Lemma~3.2 and Theorem~3.3). Chronic boredom is either a malfunction or a correct signal that cannot be acted on (its Corollary~3.4). These are claims about a psychological state, and psychology has studied the state. Over the last four decades the study of boredom has acquired validated measures, a working definition, experimental inductions and a handful of theoretical models.

This manuscript assesses what that work shows about the two theses. It does not add evidence; it sorts the evidence there is. Section~\ref{s:measures} describes the measures and definitions. Section~\ref{s:mismatch} and Section~\ref{s:function} take the two theses in turn. Section~\ref{s:attention} considers the attentional evidence for the representational claim, Section~\ref{s:induction} the limits of experimental inductions, and Section~\ref{s:test} the designs that would test the function thesis directly. Section~\ref{s:obj} answers objections to the use we make of the studies, and Section~\ref{s:not} states what the survey does not establish. The appendix gives the annotated bibliography and the computed cross-tabulations; its preface states the criteria by which studies were selected and the rules by which their bearing is computed.

\section{Measures and definitions}\label{s:measures}

\paragraph{Trait measures.} The first standard instrument was the Boredom Proneness Scale of \citet{FarmerSundberg1986}, a set of true--false items measuring the disposition to experience boredom (Far86 in the appendix). Factor analyses found it multidimensional, with factors including a need for external stimulation and a difficulty in entertaining oneself \citep{VodanovichKass1990}, and reviews have repeatedly called for further validation \citep{Vodanovich2003,VodanovichWatt2016}. \citet{Struk2017} developed a short unidimensional form (Str17). Sensation-seeking research had earlier included a boredom-susceptibility subscale \citep{Zuckerman1979}. Trait measures bear directly on Corollary~3.4 of EXI-10-a, since they measure the proneness to chronic boredom, and only indirectly on the theses about the state.

\paragraph{State measures.} The Multidimensional State Boredom Scale \citep{Fahlman2013} measures the state, with factors for disengagement, high and low arousal, inattention and time perception (Fah13). Experience-sampling studies use momentary reports \citep{LarsonRichards1991,Goetz2014,Chin2017}. State measures are what the two theses require, since both concern what the state represents and what it does.

\paragraph{Trait and state.} The distinction matters for how evidence is read. A trait correlate shows that people who are often bored differ in some respect; it does not show what an episode of boredom represents or does. Of the two theses, both concern episodes, and the evidence that bears on them most directly comes from state measures, momentary reports and inductions. About a third of the empirical entries in the appendix use trait measures, which is one reason why so many are neutral with respect to both theses: they bear instead on Corollary~3.4 of EXI-10-a, which concerns who is prone to chronic boredom and why. The cross-tabulation by design in the appendix shows the same pattern from the side of method: with respect to the function thesis, correlational and psychometric entries are neutral far more often than experiments are.

\paragraph{Definitions.} \citet{EastwoodEtAl2012} define boredom as the aversive experience of wanting, but being unable, to engage in satisfying activity, and locate the failure in attention. EXI-10-a adopts a definition in the same spirit (its Definition~2.1). The MAC model of \citet{WestgateWilson2018} separates an attentional component, a mismatch between cognitive demands and available resources in either direction, from a meaning component, the activity's failure to matter to the subject. The Boredom Feedback Model of \citet{Tam2021} treats boredom as feedback that attention is not adequately engaged. These are theoretical works; their bearing on the theses is a bearing of models on models, and the appendix records them separately from empirical studies in every cross-tabulation in which the difference matters.

\section{The mismatch thesis}\label{s:mismatch}

The mismatch thesis has three clauses, and the evidence bears on them unequally.

\paragraph{Under-challenge.} That monotony and under-stimulation produce an aversive state is the oldest finding in the area. Participants in the sensory-variation studies of \citet{Bexton1954} found prolonged low stimulation so aversive that many left despite the pay (Bex54); workplace research on repetitive tasks reached a similar conclusion \citep{OHanlon1981}; and the flow research places boredom where skill exceeds challenge \citep{Csikszentmihalyi1975,Csikszentmihalyi1990}. \citet{vanTilburgIgou2012} found that boredom is distinguished from other negative emotions by a perceived lack of challenge (Til12). The clause is well supported.

\paragraph{Meaninglessness.} \citet{Fahlman2009} found that a lack of meaning in life predicted boredom cross-sectionally, over time and under experimental manipulation (Fah09); \citet{vanTilburgIgou2012,vanTilburgIgou2017} found perceived meaninglessness to be part of the profile that distinguishes boredom from other negative emotions; and in achievement settings boredom is related to the low perceived value of the activity \citep{Pekrun2010}. The clause is the best supported by empirical studies in the list (four of its five supporting entries are empirical).

\paragraph{Over-challenge.} The third clause is the weakest. In the bibliography the only entry that supports it is a theoretical model, the MAC model, which allows the mismatch between demands and resources to run in either direction \citep{WestgateWilson2018}; no empirical study in the list isolates over-challenge as a cause of boredom. Against it stands the flow model, in both of its statements (Csi75, Csi90), on which demands exceeding skill produce anxiety rather than boredom. The computed table of clauses (Table~\ref{tab:clauses}) records \stoverpos{} supporting entry, \stoverposemp{} of them empirical, and \stoverneg{} challenging. EXI-10-a states the clause with a qualification, \enquote{over-challenge without traction}, and the qualification matters: an over-demanding task on which one makes no progress leaves attention unoccupied in the way an easy task does, while an over-demanding task one can still engage with is a source of anxiety. The flow model's categories concern the second kind. The qualification is a reasonable reading of the evidence, but the evidence does not yet test it directly, and we record the clause as supported by theory and challenged, in an unqualified form, by the flow model.

\paragraph{Fit and engagement.} The converse of the mismatch thesis, that matched challenge and skill engage, is well supported: experience sampling finds the quality of experience related to the balance of challenge and skill \citep{CsikszentmihalyiLeFevre1989,MonetaCsikszentmihalyi1996}, flow is related to the balance in field data \citep{EngeserRheinberg2008}, and an experimental manipulation of demand relative to skill produced flow, with a distinct physiological profile, in the matched condition \citep{Keller2011}. These studies (Csi89, Mon96, Eng08, Kel11) also bear on Assumption~2.4 of EXI-10-a, the value of engagement.

\paragraph{Person and situation.} A recurring finding is that individuals differ in proneness: students often bored in class were often bored elsewhere \citep{LarsonRichards1991}, and proneness is associated with attentional failures \citep{Malkovsky2012}, difficulties in self-regulation \citep{Struk2016,Mugon2018,Isacescu2017}, and alexithymia \citep{Eastwood2007}. Such findings might seem to tell against a thesis that locates boredom in a relation between person and activity. They do not, for two reasons. First, the mismatch is relational: capacities are on one side of it, and individual differences in capacity, including the capacity to sustain attention and to find activities meaningful, predict boredom on any relational view. \citet{MercerLynn2014} found that trait proneness and the situation both contributed to state boredom (Mer14), which is the pattern the relational view predicts. Second, Corollary~3.4 of EXI-10-a classifies chronic boredom as a malfunction or as a correct signal that cannot be acted on, and trait correlates are evidence about which individuals are prone to the malfunction. We code such findings as neutral with respect to the mismatch thesis and as accommodated by the corollary; the appendix states the rule and shows what changes under a stricter one.

\section{The function thesis}\label{s:function}

The function thesis has three components: boredom is a signal (Lemma~3.1 of EXI-10-a), it motivates a search for alternatives (Lemma~3.2), and the search tends to improve the agent's situation, which, with the history of selection, makes it an adaptive function (Theorem~3.3).

\paragraph{Signal.} That boredom is a distinct state, with a profile separable from apathy, anhedonia and depression \citep{Goldberg2011}, from sadness in its physiology \citep{MerrifieldDanckert2014}, and from other negative emotions in its appraisals \citep{vanTilburgIgou2017}, is a precondition of its being a dedicated signal. The theoretical models of \citet{Elpidorou2014,Elpidorou2018}, \citet{DanckertMugon2018}, \citet{DanckertEastwood2020} and \citet{Tam2021} treat it as one.

\paragraph{Motivation.} The motivational component is the best supported part of either thesis. Bored participants chose novel experiences, even unpleasant ones \citep{BenchLench2019}; they reported a wish to pursue a different goal \citep{vanTilburgIgou2012,vanTilburgIgou2017}; they changed the way they performed a repetitive task when they had to persist with it \citep{Sansone1992}, which is the second remedy Lemma~3.2 of EXI-10-a allows, a change in how one engages with the present activity; they turned to meaning-restoring responses such as nostalgia \citep{vanTilburgIgouSedikides2013} and in-group identification \citep{vanTilburgIgou2011}; and after a boring task they produced more ideas on a divergent-thinking task \citep{MannCadman2014}. The theoretical proposal of \citet{BenchLench2013}, that boredom prompts the pursuit of alternative goals, is the empirical counterpart of Lemma~3.2.

\paragraph{The shock studies.} The most widely cited finding is that of \citet{WilsonEtAl2014} (Wil14): participants generally did not enjoy several minutes alone with their thoughts, and many gave themselves mild electric shocks rather than sit without stimulation. Laboratory boredom inductions have since been found to increase eating and self-administered shocks \citep{Havermans2015} and self-administered painful stimulation \citep{Nederkoorn2016} (Hav15, Ned16). EXI-10-a cites the first study as evidence of the motivational profile. The rules of the appendix code these studies as bearing in both directions on the function thesis: they show the motivation to escape monotony, and they show that the motivation can be satisfied by responses that are costly. The reply of EXI-10-a to its second objection applies: boredom motivates a search for engagement but does not select the options on offer, and in a laboratory in which the only available stimulation is a shock, the search ends there. The studies show that the signal is strong enough to override a mild aversion to pain, which is what an effective signal requires; they do not show that the responses it produces are adaptive in the environments in which the disposition was retained, and no study could show this directly.

\paragraph{Benefit.} The adaptive component is the least directly tested. The evidence that engagement yields concentration, involvement and better experience (Section~\ref{s:mismatch}) supports Assumption~2.4; the evidence that boredom prompts goal change supports Lemma~3.2; and the conjunction of the two is the argument of Theorem~3.3. There is no study in the list that follows bored agents to measure whether the activities they turn to are better matched to their capacities than those they leave, which is the direct test of the claim that the search tends to reduce mismatch. The negative relation between boredom and academic achievement \citep{Pekrun2010,Tze2016} and the association of trait proneness with health symptoms \citep{SommersVodanovich2000} and with rule-breaking during the pandemic \citep{Wolff2020,Boylan2021} concern costs of boredom within situations one cannot leave, or of chronic proneness; we code them as accommodated by Corollary~3.4 and show in the appendix how the counts change if they are instead counted as challenges.

\paragraph{Summary.} Of the \stn{} entries, \staS{} support the mismatch thesis, \staSandC{} bear on it in both directions, and \staN{} are neutral; \stbS{} support the function thesis, \stbSandC{} bear in both directions, and \stbN{} are neutral. Under the strict rule for costs, \stbsC{} entries challenge the function thesis outright and \stbsSandC{} bear in both directions. No entry challenges either thesis outright under the lenient rule.


\section{Attention and the content of boredom}\label{s:attention}

Lemma~3.1 of EXI-10-a says that boredom represents a mismatch and is correct when there is one. The empirical counterpart of a representational claim is a claim about what reliably triggers the state and what the state carries information about. The attentional literature bears on this most directly. \citet{OHanlon1981} described boredom in repetitive work as the felt cost of sustaining attention to a task that does not hold it. \citet{Malkovsky2012} found boredom proneness associated with failures of sustained attention (Mal12), and \citet{HunterEastwood2018} examined how trait boredom, state boredom and sustained attention are related (Hun18). \citet{DanckertMerrifield2018} associated boredom induction with activity in the default mode network, the network more active when attention is not directed to an external task (DnM18), and \citet{Danckert2018restless} found boredom to combine low arousal with restlessness (DnH18). \citet{EastwoodEtAl2012} built these findings into their definition, and \citet{Tam2021} into the claim that boredom is feedback about attentional engagement.

Two cautions apply. First, the attentional findings fit the mismatch thesis on its attentional reading, on which the capacities in question include the capacity to attend; they say less about its meaning clause, which concerns the capacity to care. The MAC model treats the two as separate routes to the same state \citep{WestgateWilson2018}, and EXI-10-a's Definition~2.2 includes both as forms of a single mismatch. The empirical literature has so far studied them largely separately: the entries coded \texttt{attn} and those coded \texttt{meaning+} overlap only in the MAC model itself (Table~\ref{tab:clauses}). Second, the direction of explanation is not settled. If attentional failure causes boredom, boredom represents the failure; if boredom causes attentional failure, the correlation shows a cost of the state rather than its content. Hunter and Eastwood's question, whether state boredom causes failures of attention, is posed in exactly these terms. For Lemma~3.1 the first direction is required, and the correlational studies are consistent with it without establishing it.

\section{Inductions and their limits}\label{s:induction}

Most experimental evidence about the function thesis comes from inductions: participants copy references, watch a repetitive film, turn pegs or sit alone, and their subsequent choices are compared with those of participants in a comparison condition. The method establishes that the induced state has the motivational effects measured (Wil14, Hav15, Ned16, Ben19, Man14), but it has two limits that matter for the thesis. The first is that the alternatives on offer are fixed by the experimenter and are usually few: a shock button, a bowl of food, a set of images, a creativity task. The function thesis is that boredom moves the agent toward activities that match her capacities, and an induction that offers no such activity cannot show this, though it can show the strength and direction of the push. The second is that inductions produce boredom by under-stimulation. None of the experimental entries in the list induces boredom by over-challenge or by meaninglessness alone, except the meaning manipulation reported by \citet{Fahlman2009}; the motivational evidence therefore concerns one of the three forms of mismatch. The flow experiment of \citet{Keller2011} (Kel11) varies demand in both directions, but it measures flow and its physiology in each condition, not the choices that follow; it shows what a match does, not what boredom prompts.

\section{What would test the function thesis}\label{s:test}

The gap between the motivational evidence and the adaptive claim can be stated as a prediction that existing designs do not test: when bored agents are free to choose among activities that differ in how well they match their capacities, boredom should increase the probability of a move to a better-matched activity, and the move should be followed by a decline in boredom and an increase in engagement or learning. Three designs would test it. An experiment could offer, after a boredom induction, a menu of tasks graded in difficulty relative to a skill measured in advance, and record which is chosen; the thesis predicts that bored participants move toward the matched level more than controls. An experience-sampling study could record momentary boredom, the activity, and its subjective challenge and skill, and test whether boredom predicts a change of activity toward balance at the next sampling point; the experience-sampling studies in the list (Csi89, Mon96, Lar91, Goe14, Chi17) record the ingredients but do not report this transition analysis. A longitudinal study of learners could test whether the episodes of boredom that are followed by a change of task or approach are those after which performance improves. The selected-effect history that Definition~2.3 of EXI-10-a requires would not be tested by any of these, but its present-day counterpart, that the disposition has the effect for which it would have been selected, would be.

\section{Relations to other families}

The survey bears on three neighboring conclusions. EXI-17 holds that a meaningful life requires active engagement with something worthwhile; the meaning clause of the mismatch thesis is the affective counterpart of that view, and the evidence that lack of meaning predicts boredom (Fah09, Til12, Til17) is evidence that the mind registers the absence of such engagement. EXI-03 holds that an endless life would exhaust the categorical desires that give life its point; EXI-10-a draws the consequence that boredom would then become a correct signal that nothing could answer, and the trait findings of Corollary~3.4 are the nearest empirical analogue, a correct or incorrect signal that persists. PHE-04 holds that profound boredom, in Heidegger's sense, is intentional toward time itself; the empirical literature does not study that state, and the time-estimation findings of \citet{DanckertAllman2005} concern the ordinary state's distortion of experienced duration, not the profound state's object.

\section{Objections and replies}\label{s:obj}

\begin{objection}[Selection]
A bibliography assembled by advocates of a thesis will favor studies that support it.
\end{objection}

\begin{reply}
The criteria of selection, stated in the preface to the appendix, are confidence in the bibliographic record and relevance to boredom as a state or a trait, not the direction of the finding; the list includes every study we could verify on the costs of boredom and on trait correlates, which the lenient rule codes as neutral and the strict rule as challenges. The list is not exhaustive, and a reader who knows of a study that challenges either thesis can code it by the same rules. What the selection cannot do is turn a correlational literature into a test of a selected-effect claim; Section~\ref{s:not} says so.
\end{reply}

\begin{objection}[Coding is interpretation]
The codes assigned to each study are judgments, and the computed bearing inherits them.
\end{objection}

\begin{reply}
That is correct, and the appendix makes the dependence explicit: each entry shows its codes and the trace of the rule applied to them, so that a disagreement about a bearing can be located in a disagreement about a code. The rules themselves are fixed in advance and applied uniformly. The codes record what a study found, not whether it supports a thesis; the rules convert one into the other.
\end{reply}

\begin{objection}[Correlates of proneness are counterevidence]
If boredom proneness predicts depression, poor health and rule-breaking, a disposition to boredom is harmful, and the claim that it is adaptive is false.
\end{objection}

\begin{reply}
The function thesis concerns the disposition of normal agents to become bored when there is a mismatch, not the high end of a distribution of proneness. That some individuals are bored too often, or in situations they cannot change, is predicted by Corollary~3.4 of EXI-10-a and does not tell against the function of the normal disposition, any more than chronic pain tells against the function of pain. The strict-rule column of the appendix shows how much of the picture changes if one refuses this distinction: \stbsC{} entries become challenges, and the motivational evidence is unaffected.
\end{reply}

\section{What the survey does not establish}\label{s:not}

The survey does not establish the selected-effect history that Definition~2.3 of EXI-10-a requires. Function in that sense is a matter of why the disposition was retained, and the studies are synchronic: they show what boredom does now, in laboratories, schools and everyday life. They are evidence for the function claim only through the inference from present effects to past selection, which EXI-10-a makes and which the studies neither confirm nor refute. The survey does not establish the over-challenge clause, for the reason given in Section~\ref{s:mismatch}. And it does not establish that the list is complete: it is limited to studies whose details we could state with confidence.

\section{Conclusion}

The empirical study of boredom supports the mismatch thesis in its under-challenge and meaning clauses and in its converse, that matched engagement absorbs, and leaves its over-challenge clause resting on theory and contested by the flow model. It supports the function thesis chiefly through its motivational component: bored people seek alternatives, change how they engage, and turn to sources of meaning. The adaptive component remains an inference from these findings together with the value of engagement, and the studies of costly responses and of chronic proneness are consistent with it once the function of the normal disposition is distinguished from its failures.

\appendix

\section{Annotated bibliography}\label{app:a}

\paragraph{Selection criteria.} An entry was included only if (i) we are confident that the work exists with the stated authors, year, title and venue; (ii) it reports an empirical study of boredom as a state or a trait, or of the challenge--skill balance on which flow research bears, or it is a theoretical model or review that the empirical literature uses as a framework; and (iii) its main finding can be stated without figures we could not verify. Page ranges are omitted where we were not certain of them. Sample sizes and effect sizes are not reported; findings are stated in qualitative terms and cautiously. The direction of a study's finding played no role in selection. The list is not a systematic review and makes no claim to exhaustiveness.

\paragraph{Encoding.} Each entry has an identifier (author abbreviation and year), a design (experiment, experience sampling, correlational, psychometric, physiological or imaging, meta-analysis, review, theory, interview-based), a short description of the sample and the measure, a statement of the main finding, and a set of codes recording what the study found (Table~\ref{tab:codes}). The codes are assigned by us from the finding; they do not record a bearing.

\paragraph{Rules.} The program computes the bearing of each entry on (a) the mismatch thesis and (b) the function thesis as follows. (a) An entry supports the thesis if it carries one of \texttt{under+}, \texttt{over+}, \texttt{meaning+}, \texttt{fit-}, \texttt{attn} or \texttt{relational}, and challenges it if it carries \texttt{overanx}. (b) An entry supports the thesis if it carries a motivational code (\texttt{seek}, \texttt{switch}, \texttt{meaningreg}, \texttt{creat}) or a signal code (\texttt{signal}, \texttt{distinct}), and challenges it if it carries \texttt{harm}, a costly response chosen to escape boredom. In each case an entry that both supports and challenges is recorded as S/C rather than netted out, and an entry that does neither as N. The codes \texttt{chronic}, \texttt{perf-} and \texttt{person} are recorded as accommodated by Corollary~3.4 of EXI-10-a and count as neutral; under the \emph{strict} rule, \texttt{chronic} and \texttt{perf-} count as challenges to (b). Both columns are computed and tabulated.

\paragraph{Summary.} There are \stn{} entries, \stnemp{} empirical and \stntheo{} reviews or theoretical works, published between \stymin{} and \stymax{}. For (a): \staS{} S, \staSandC{} S/C, \staN{} N, \staC{} C. For (b), lenient rule: \stbS{} S, \stbSandC{} S/C, \stbN{} N, \stbC{} C; strict rule: \stbsS{} S, \stbsSandC{} S/C, \stbsN{} N, \stbsC{} C.

\needspace{6\baselineskip}
\noindent\textbf{Bex54}\enspace \citet{Bexton1954}.\par
{\small\noindent\textit{Design:} experiment.\quad \textit{Sample:} paid male students.\quad \textit{Measure:} tolerance of reduced sensory variation; cognitive tests.\par
\noindent\textit{Finding.} Prolonged confinement with little sensory variation was found highly aversive; many participants ended their participation early despite the pay, and some cognitive performance declined.\par
\noindent\textit{Codes:} \texttt{under+ seek}.\quad \textit{Bearing:} (a) \textbf{S}, (b) \textbf{S}.\par
\noindent\textit{Rule trace:} (a) S: support {[}under+{]}, challenge {[}--{]}; (b) S: motivation {[}seek{]}, signal {[}--{]}, cost {[}--{]}.\par}
\medskip
\needspace{6\baselineskip}
\noindent\textbf{Ber60}\enspace \citet{Berlyne1960}.\par
{\small\noindent\textit{Design:} theory.\quad \textit{Sample:} none (theory and review).\quad \textit{Measure:} --.\par
\noindent\textit{Finding.} Proposes that organisms seek an intermediate level of arousal and that monotony motivates exploratory and diversive behavior.\par
\noindent\textit{Codes:} \texttt{under+ seek}.\quad \textit{Bearing:} (a) \textbf{S}, (b) \textbf{S}.\par
\noindent\textit{Rule trace:} (a) S: support {[}under+{]}, challenge {[}--{]}; (b) S: motivation {[}seek{]}, signal {[}--{]}, cost {[}--{]}.\par}
\medskip
\needspace{6\baselineskip}
\noindent\textbf{Csi75}\enspace \citet{Csikszentmihalyi1975}.\par
{\small\noindent\textit{Design:} interview-based.\quad \textit{Sample:} adults engaged in demanding activities.\quad \textit{Measure:} interviews and questionnaires.\par
\noindent\textit{Finding.} Introduces the flow model: absorption where challenge and skill are matched, boredom where skill exceeds challenge, and anxiety, rather than boredom, where challenge exceeds skill.\par
\noindent\textit{Codes:} \texttt{under+ fit- overanx}.\quad \textit{Bearing:} (a) \textbf{S/C}, (b) \textbf{N}.\par
\noindent\textit{Rule trace:} (a) S/C: support {[}fit-, under+{]}, challenge {[}overanx{]}; (b) N: motivation {[}--{]}, signal {[}--{]}, cost {[}--{]}.\par}
\medskip
\needspace{6\baselineskip}
\noindent\textbf{Zuc79}\enspace \citet{Zuckerman1979}.\par
{\small\noindent\textit{Design:} psychometric.\quad \textit{Sample:} several adult samples.\quad \textit{Measure:} Sensation Seeking Scale (Boredom Susceptibility subscale).\par
\noindent\textit{Finding.} Presents sensation seeking as a trait with a boredom-susceptibility component: individuals differ in their aversion to repetition and routine.\par
\noindent\textit{Codes:} \texttt{meas person}.\quad \textit{Bearing:} (a) \textbf{N}, (b) \textbf{N}.\par
\noindent\textit{Rule trace:} (a) N: support {[}--{]}, challenge {[}--{]}; (b) N: motivation {[}--{]}, signal {[}--{]}, cost {[}--{]}; accommodated by Cor.~3.4 of EXI-10-a: {[}person{]}.\par}
\medskip
\needspace{6\baselineskip}
\noindent\textbf{OHa81}\enspace \citet{OHanlon1981}.\par
{\small\noindent\textit{Design:} review.\quad \textit{Sample:} none (review of workplace studies).\quad \textit{Measure:} --.\par
\noindent\textit{Finding.} Reviews boredom in repetitive work and proposes that it arises from sustained effort to attend to monotonous tasks.\par
\noindent\textit{Codes:} \texttt{under+ attn}.\quad \textit{Bearing:} (a) \textbf{S}, (b) \textbf{N}.\par
\noindent\textit{Rule trace:} (a) S: support {[}attn, under+{]}, challenge {[}--{]}; (b) N: motivation {[}--{]}, signal {[}--{]}, cost {[}--{]}.\par}
\medskip
\needspace{6\baselineskip}
\noindent\textbf{Far86}\enspace \citet{FarmerSundberg1986}.\par
{\small\noindent\textit{Design:} psychometric.\quad \textit{Sample:} undergraduates.\quad \textit{Measure:} Boredom Proneness Scale (BPS).\par
\noindent\textit{Finding.} Develops a 28-item true--false scale of the disposition to experience boredom, with reported reliability and correlations with measures of depression, hopelessness and loneliness.\par
\noindent\textit{Codes:} \texttt{meas chronic}.\quad \textit{Bearing:} (a) \textbf{N}, (b) \textbf{N}.\par
\noindent\textit{Rule trace:} (a) N: support {[}--{]}, challenge {[}--{]}; (b) N: motivation {[}--{]}, signal {[}--{]}, cost {[}--{]}; accommodated by Cor.~3.4 of EXI-10-a: {[}chronic{]}.\par}
\medskip
\needspace{6\baselineskip}
\noindent\textbf{Csi89}\enspace \citet{CsikszentmihalyiLeFevre1989}.\par
{\small\noindent\textit{Design:} experience sampling.\quad \textit{Sample:} adult workers.\quad \textit{Measure:} experience sampling of challenge, skill and quality of experience.\par
\noindent\textit{Finding.} Moments of high and balanced challenge and skill were reported more often at work than in leisure and were associated with more positive experience.\par
\noindent\textit{Codes:} \texttt{fit-}.\quad \textit{Bearing:} (a) \textbf{S}, (b) \textbf{N}.\par
\noindent\textit{Rule trace:} (a) S: support {[}fit-{]}, challenge {[}--{]}; (b) N: motivation {[}--{]}, signal {[}--{]}, cost {[}--{]}.\par}
\medskip
\needspace{6\baselineskip}
\noindent\textbf{Csi90}\enspace \citet{Csikszentmihalyi1990}.\par
{\small\noindent\textit{Design:} theory.\quad \textit{Sample:} none (synthesis).\quad \textit{Measure:} --.\par
\noindent\textit{Finding.} Synthesizes flow research: optimal experience where challenge matches skill; boredom where skill exceeds challenge; anxiety where challenge exceeds skill.\par
\noindent\textit{Codes:} \texttt{under+ fit- overanx}.\quad \textit{Bearing:} (a) \textbf{S/C}, (b) \textbf{N}.\par
\noindent\textit{Rule trace:} (a) S/C: support {[}fit-, under+{]}, challenge {[}overanx{]}; (b) N: motivation {[}--{]}, signal {[}--{]}, cost {[}--{]}.\par}
\medskip
\needspace{6\baselineskip}
\noindent\textbf{Vod90}\enspace \citet{VodanovichKass1990}.\par
{\small\noindent\textit{Design:} psychometric.\quad \textit{Sample:} undergraduates.\quad \textit{Measure:} BPS (factor analysis).\par
\noindent\textit{Finding.} Finds that the BPS is multidimensional, with factors including the need for external stimulation and difficulty in self-entertainment.\par
\noindent\textit{Codes:} \texttt{meas}.\quad \textit{Bearing:} (a) \textbf{N}, (b) \textbf{N}.\par
\noindent\textit{Rule trace:} (a) N: support {[}--{]}, challenge {[}--{]}; (b) N: motivation {[}--{]}, signal {[}--{]}, cost {[}--{]}.\par}
\medskip
\needspace{6\baselineskip}
\noindent\textbf{Lar91}\enspace \citet{LarsonRichards1991}.\par
{\small\noindent\textit{Design:} experience sampling.\quad \textit{Sample:} middle-school students.\quad \textit{Measure:} experience sampling of boredom.\par
\noindent\textit{Finding.} Boredom was frequent both in and out of school, and students who were often bored in class tended to be often bored elsewhere.\par
\noindent\textit{Codes:} \texttt{person}.\quad \textit{Bearing:} (a) \textbf{N}, (b) \textbf{N}.\par
\noindent\textit{Rule trace:} (a) N: support {[}--{]}, challenge {[}--{]}; (b) N: motivation {[}--{]}, signal {[}--{]}, cost {[}--{]}; accommodated by Cor.~3.4 of EXI-10-a: {[}person{]}.\par}
\medskip
\needspace{6\baselineskip}
\noindent\textbf{San92}\enspace \citet{Sansone1992}.\par
{\small\noindent\textit{Design:} experiment.\quad \textit{Sample:} undergraduates.\quad \textit{Measure:} interest and task-variation strategies on a repetitive task.\par
\noindent\textit{Finding.} When given a reason to persist at a boring task, participants varied the way they performed it, making it more interesting.\par
\noindent\textit{Codes:} \texttt{switch}.\quad \textit{Bearing:} (a) \textbf{N}, (b) \textbf{S}.\par
\noindent\textit{Rule trace:} (a) N: support {[}--{]}, challenge {[}--{]}; (b) S: motivation {[}switch{]}, signal {[}--{]}, cost {[}--{]}.\par}
\medskip
\needspace{6\baselineskip}
\noindent\textbf{Mon96}\enspace \citet{MonetaCsikszentmihalyi1996}.\par
{\small\noindent\textit{Design:} experience sampling.\quad \textit{Sample:} high-school students.\quad \textit{Measure:} experience sampling of challenge, skill and experience.\par
\noindent\textit{Finding.} Perceived challenge, skill and their balance were related to concentration, involvement and other aspects of the quality of experience.\par
\noindent\textit{Codes:} \texttt{fit-}.\quad \textit{Bearing:} (a) \textbf{S}, (b) \textbf{N}.\par
\noindent\textit{Rule trace:} (a) S: support {[}fit-{]}, challenge {[}--{]}; (b) N: motivation {[}--{]}, signal {[}--{]}, cost {[}--{]}.\par}
\medskip
\needspace{6\baselineskip}
\noindent\textbf{Som00}\enspace \citet{SommersVodanovich2000}.\par
{\small\noindent\textit{Design:} correlational.\quad \textit{Sample:} undergraduates.\quad \textit{Measure:} BPS; symptom checklists.\par
\noindent\textit{Finding.} Boredom proneness was associated with reports of psychological and physical-health symptoms.\par
\noindent\textit{Codes:} \texttt{chronic}.\quad \textit{Bearing:} (a) \textbf{N}, (b) \textbf{N}.\par
\noindent\textit{Rule trace:} (a) N: support {[}--{]}, challenge {[}--{]}; (b) N: motivation {[}--{]}, signal {[}--{]}, cost {[}--{]}; accommodated by Cor.~3.4 of EXI-10-a: {[}chronic{]}.\par}
\medskip
\needspace{6\baselineskip}
\noindent\textbf{Vod03}\enspace \citet{Vodanovich2003}.\par
{\small\noindent\textit{Design:} review.\quad \textit{Sample:} none (review).\quad \textit{Measure:} review of boredom scales.\par
\noindent\textit{Finding.} Reviews the psychometric properties of trait and state boredom measures and recommends further validation work.\par
\noindent\textit{Codes:} \texttt{meas}.\quad \textit{Bearing:} (a) \textbf{N}, (b) \textbf{N}.\par
\noindent\textit{Rule trace:} (a) N: support {[}--{]}, challenge {[}--{]}; (b) N: motivation {[}--{]}, signal {[}--{]}, cost {[}--{]}.\par}
\medskip
\needspace{6\baselineskip}
\noindent\textbf{Dan05}\enspace \citet{DanckertAllman2005}.\par
{\small\noindent\textit{Design:} correlational.\quad \textit{Sample:} healthy adults high or low in boredom proneness.\quad \textit{Measure:} time-estimation and attention tasks; boredom proneness.\par
\noindent\textit{Finding.} High and low boredom-prone participants did not differ in the temporal allocation of attention, but high boredom-prone participants made larger errors in estimating temporal intervals.\par
\noindent\textit{Codes:} \texttt{distinct}.\quad \textit{Bearing:} (a) \textbf{N}, (b) \textbf{S}.\par
\noindent\textit{Rule trace:} (a) N: support {[}--{]}, challenge {[}--{]}; (b) S: motivation {[}--{]}, signal {[}distinct{]}, cost {[}--{]}.\par}
\medskip
\needspace{6\baselineskip}
\noindent\textbf{Eas07}\enspace \citet{Eastwood2007}.\par
{\small\noindent\textit{Design:} correlational.\quad \textit{Sample:} undergraduates.\quad \textit{Measure:} BPS; alexithymia scale.\par
\noindent\textit{Finding.} Boredom proneness was associated with difficulty in identifying and describing one's feelings.\par
\noindent\textit{Codes:} \texttt{person}.\quad \textit{Bearing:} (a) \textbf{N}, (b) \textbf{N}.\par
\noindent\textit{Rule trace:} (a) N: support {[}--{]}, challenge {[}--{]}; (b) N: motivation {[}--{]}, signal {[}--{]}, cost {[}--{]}; accommodated by Cor.~3.4 of EXI-10-a: {[}person{]}.\par}
\medskip
\needspace{6\baselineskip}
\noindent\textbf{Eng08}\enspace \citet{EngeserRheinberg2008}.\par
{\small\noindent\textit{Design:} correlational.\quad \textit{Sample:} students and adults in learning and leisure settings.\quad \textit{Measure:} flow and challenge--skill ratings.\par
\noindent\textit{Finding.} Flow was related to the balance of challenge and skill, with the relation moderated by individual and situational factors.\par
\noindent\textit{Codes:} \texttt{fit-}.\quad \textit{Bearing:} (a) \textbf{S}, (b) \textbf{N}.\par
\noindent\textit{Rule trace:} (a) S: support {[}fit-{]}, challenge {[}--{]}; (b) N: motivation {[}--{]}, signal {[}--{]}, cost {[}--{]}.\par}
\medskip
\needspace{6\baselineskip}
\noindent\textbf{Fah09}\enspace \citet{Fahlman2009}.\par
{\small\noindent\textit{Design:} correlational.\quad \textit{Sample:} undergraduates (cross-sectional, longitudinal and experimental studies).\quad \textit{Measure:} BPS; meaning-in-life measures.\par
\noindent\textit{Finding.} Lack of meaning in life predicted boredom, including over time and, in an experimental study, when meaning was manipulated.\par
\noindent\textit{Codes:} \texttt{meaning+}.\quad \textit{Bearing:} (a) \textbf{S}, (b) \textbf{N}.\par
\noindent\textit{Rule trace:} (a) S: support {[}meaning+{]}, challenge {[}--{]}; (b) N: motivation {[}--{]}, signal {[}--{]}, cost {[}--{]}.\par}
\medskip
\needspace{6\baselineskip}
\noindent\textbf{Pek10}\enspace \citet{Pekrun2010}.\par
{\small\noindent\textit{Design:} correlational.\quad \textit{Sample:} university students.\quad \textit{Measure:} achievement-emotion questionnaire scales.\par
\noindent\textit{Finding.} Boredom in achievement settings was related to appraisals of control and of the value of the activity, and predicted lower achievement.\par
\noindent\textit{Codes:} \texttt{meaning+ perf-}.\quad \textit{Bearing:} (a) \textbf{S}, (b) \textbf{N}.\par
\noindent\textit{Rule trace:} (a) S: support {[}meaning+{]}, challenge {[}--{]}; (b) N: motivation {[}--{]}, signal {[}--{]}, cost {[}--{]}; accommodated by Cor.~3.4 of EXI-10-a: {[}perf-{]}.\par}
\medskip
\needspace{6\baselineskip}
\noindent\textbf{Gol11}\enspace \citet{Goldberg2011}.\par
{\small\noindent\textit{Design:} psychometric.\quad \textit{Sample:} undergraduates.\quad \textit{Measure:} boredom, apathy, anhedonia and depression scales.\par
\noindent\textit{Finding.} Boredom was related to, but separable from, apathy, anhedonia and depression.\par
\noindent\textit{Codes:} \texttt{distinct}.\quad \textit{Bearing:} (a) \textbf{N}, (b) \textbf{S}.\par
\noindent\textit{Rule trace:} (a) N: support {[}--{]}, challenge {[}--{]}; (b) S: motivation {[}--{]}, signal {[}distinct{]}, cost {[}--{]}.\par}
\medskip
\needspace{6\baselineskip}
\noindent\textbf{Kel11}\enspace \citet{Keller2011}.\par
{\small\noindent\textit{Design:} experiment.\quad \textit{Sample:} laboratory sample.\quad \textit{Measure:} flow ratings; heart-rate variability; salivary cortisol.\par
\noindent\textit{Finding.} A task whose demands were matched to the participant's skill, compared with too easy and too hard versions, produced flow and a distinct physiological profile.\par
\noindent\textit{Codes:} \texttt{fit-}.\quad \textit{Bearing:} (a) \textbf{S}, (b) \textbf{N}.\par
\noindent\textit{Rule trace:} (a) S: support {[}fit-{]}, challenge {[}--{]}; (b) N: motivation {[}--{]}, signal {[}--{]}, cost {[}--{]}.\par}
\medskip
\needspace{6\baselineskip}
\noindent\textbf{Til11}\enspace \citet{vanTilburgIgou2011}.\par
{\small\noindent\textit{Design:} experiment.\quad \textit{Sample:} student and adult samples.\quad \textit{Measure:} boredom inductions; in-group evaluation.\par
\noindent\textit{Finding.} Induced boredom increased the favoring of the in-group, interpreted as a way of re-establishing meaning.\par
\noindent\textit{Codes:} \texttt{meaningreg}.\quad \textit{Bearing:} (a) \textbf{N}, (b) \textbf{S}.\par
\noindent\textit{Rule trace:} (a) N: support {[}--{]}, challenge {[}--{]}; (b) S: motivation {[}meaningreg{]}, signal {[}--{]}, cost {[}--{]}.\par}
\medskip
\needspace{6\baselineskip}
\noindent\textbf{Eas12}\enspace \citet{EastwoodEtAl2012}.\par
{\small\noindent\textit{Design:} theory.\quad \textit{Sample:} none (theory).\quad \textit{Measure:} --.\par
\noindent\textit{Finding.} Defines boredom as the aversive experience of wanting, but being unable, to engage in satisfying activity, and grounds it in failures of attention.\par
\noindent\textit{Codes:} \texttt{attn signal}.\quad \textit{Bearing:} (a) \textbf{S}, (b) \textbf{S}.\par
\noindent\textit{Rule trace:} (a) S: support {[}attn{]}, challenge {[}--{]}; (b) S: motivation {[}--{]}, signal {[}signal{]}, cost {[}--{]}.\par}
\medskip
\needspace{6\baselineskip}
\noindent\textbf{Mal12}\enspace \citet{Malkovsky2012}.\par
{\small\noindent\textit{Design:} experiment.\quad \textit{Sample:} undergraduates.\quad \textit{Measure:} BPS; sustained-attention task.\par
\noindent\textit{Finding.} Boredom proneness was associated with failures of sustained attention.\par
\noindent\textit{Codes:} \texttt{attn person}.\quad \textit{Bearing:} (a) \textbf{S}, (b) \textbf{N}.\par
\noindent\textit{Rule trace:} (a) S: support {[}attn{]}, challenge {[}--{]}; (b) N: motivation {[}--{]}, signal {[}--{]}, cost {[}--{]}; accommodated by Cor.~3.4 of EXI-10-a: {[}person{]}.\par}
\medskip
\needspace{6\baselineskip}
\noindent\textbf{Til12}\enspace \citet{vanTilburgIgou2012}.\par
{\small\noindent\textit{Design:} experiment.\quad \textit{Sample:} student and adult samples.\quad \textit{Measure:} recalled and induced boredom; appraisal ratings.\par
\noindent\textit{Finding.} Boredom was characterized, relative to other negative emotions, by perceived lack of challenge and lack of meaning, and by a wish to do something else.\par
\noindent\textit{Codes:} \texttt{under+ meaning+ signal switch}.\quad \textit{Bearing:} (a) \textbf{S}, (b) \textbf{S}.\par
\noindent\textit{Rule trace:} (a) S: support {[}meaning+, under+{]}, challenge {[}--{]}; (b) S: motivation {[}switch{]}, signal {[}signal{]}, cost {[}--{]}.\par}
\medskip
\needspace{6\baselineskip}
\noindent\textbf{Ben13}\enspace \citet{BenchLench2013}.\par
{\small\noindent\textit{Design:} theory.\quad \textit{Sample:} none (theory and review).\quad \textit{Measure:} --.\par
\noindent\textit{Finding.} Argues that boredom motivates the pursuit of new goals when the current situation has ceased to be rewarding.\par
\noindent\textit{Codes:} \texttt{switch}.\quad \textit{Bearing:} (a) \textbf{N}, (b) \textbf{S}.\par
\noindent\textit{Rule trace:} (a) N: support {[}--{]}, challenge {[}--{]}; (b) S: motivation {[}switch{]}, signal {[}--{]}, cost {[}--{]}.\par}
\medskip
\needspace{6\baselineskip}
\noindent\textbf{Fah13}\enspace \citet{Fahlman2013}.\par
{\small\noindent\textit{Design:} psychometric.\quad \textit{Sample:} undergraduates and community adults.\quad \textit{Measure:} Multidimensional State Boredom Scale (MSBS).\par
\noindent\textit{Finding.} Develops a state boredom scale with factors for disengagement, high and low arousal, inattention and time perception.\par
\noindent\textit{Codes:} \texttt{meas}.\quad \textit{Bearing:} (a) \textbf{N}, (b) \textbf{N}.\par
\noindent\textit{Rule trace:} (a) N: support {[}--{]}, challenge {[}--{]}; (b) N: motivation {[}--{]}, signal {[}--{]}, cost {[}--{]}.\par}
\medskip
\needspace{6\baselineskip}
\noindent\textbf{Til13}\enspace \citet{vanTilburgIgouSedikides2013}.\par
{\small\noindent\textit{Design:} experiment.\quad \textit{Sample:} student and adult samples.\quad \textit{Measure:} boredom inductions; nostalgia; meaning.\par
\noindent\textit{Finding.} Induced boredom increased nostalgia, which in turn increased perceived meaning.\par
\noindent\textit{Codes:} \texttt{meaningreg}.\quad \textit{Bearing:} (a) \textbf{N}, (b) \textbf{S}.\par
\noindent\textit{Rule trace:} (a) N: support {[}--{]}, challenge {[}--{]}; (b) S: motivation {[}meaningreg{]}, signal {[}--{]}, cost {[}--{]}.\par}
\medskip
\needspace{6\baselineskip}
\noindent\textbf{Elp14}\enspace \citet{Elpidorou2014}.\par
{\small\noindent\textit{Design:} theory.\quad \textit{Sample:} none (theory).\quad \textit{Measure:} --.\par
\noindent\textit{Finding.} Argues that boredom is a regulatory state that informs the agent that her situation does not fit her goals and motivates a change.\par
\noindent\textit{Codes:} \texttt{signal switch}.\quad \textit{Bearing:} (a) \textbf{N}, (b) \textbf{S}.\par
\noindent\textit{Rule trace:} (a) N: support {[}--{]}, challenge {[}--{]}; (b) S: motivation {[}switch{]}, signal {[}signal{]}, cost {[}--{]}.\par}
\medskip
\needspace{6\baselineskip}
\noindent\textbf{Ger14}\enspace \citet{Gerritsen2014}.\par
{\small\noindent\textit{Design:} correlational.\quad \textit{Sample:} undergraduates.\quad \textit{Measure:} boredom proneness; inattention, hyperactivity, impulsivity and executive-function measures.\par
\noindent\textit{Finding.} Inattention, hyperactivity and executive dysfunction predicted boredom proneness, largely through the variance they share.\par
\noindent\textit{Codes:} \texttt{person}.\quad \textit{Bearing:} (a) \textbf{N}, (b) \textbf{N}.\par
\noindent\textit{Rule trace:} (a) N: support {[}--{]}, challenge {[}--{]}; (b) N: motivation {[}--{]}, signal {[}--{]}, cost {[}--{]}; accommodated by Cor.~3.4 of EXI-10-a: {[}person{]}.\par}
\medskip
\needspace{6\baselineskip}
\noindent\textbf{Goe14}\enspace \citet{Goetz2014}.\par
{\small\noindent\textit{Design:} experience sampling.\quad \textit{Sample:} high-school and university students.\quad \textit{Measure:} experience sampling of boredom, valence and arousal.\par
\noindent\textit{Finding.} Distinguishes several types of boredom differing in valence and arousal, among them a searching type in which the subject looks for alternatives.\par
\noindent\textit{Codes:} \texttt{meas seek}.\quad \textit{Bearing:} (a) \textbf{N}, (b) \textbf{S}.\par
\noindent\textit{Rule trace:} (a) N: support {[}--{]}, challenge {[}--{]}; (b) S: motivation {[}seek{]}, signal {[}--{]}, cost {[}--{]}.\par}
\medskip
\needspace{6\baselineskip}
\noindent\textbf{Man14}\enspace \citet{MannCadman2014}.\par
{\small\noindent\textit{Design:} experiment.\quad \textit{Sample:} adult participants.\quad \textit{Measure:} boredom induction; divergent-thinking task.\par
\noindent\textit{Finding.} Participants who first performed a boring task produced more ideas on a subsequent divergent-thinking task than controls.\par
\noindent\textit{Codes:} \texttt{creat}.\quad \textit{Bearing:} (a) \textbf{N}, (b) \textbf{S}.\par
\noindent\textit{Rule trace:} (a) N: support {[}--{]}, challenge {[}--{]}; (b) S: motivation {[}creat{]}, signal {[}--{]}, cost {[}--{]}.\par}
\medskip
\needspace{6\baselineskip}
\noindent\textbf{Mer14}\enspace \citet{MercerLynn2014}.\par
{\small\noindent\textit{Design:} correlational.\quad \textit{Sample:} undergraduates.\quad \textit{Measure:} trait and state boredom; situational ratings.\par
\noindent\textit{Finding.} Both trait boredom proneness and features of the situation contributed to state boredom.\par
\noindent\textit{Codes:} \texttt{relational person}.\quad \textit{Bearing:} (a) \textbf{S}, (b) \textbf{N}.\par
\noindent\textit{Rule trace:} (a) S: support {[}relational{]}, challenge {[}--{]}; (b) N: motivation {[}--{]}, signal {[}--{]}, cost {[}--{]}; accommodated by Cor.~3.4 of EXI-10-a: {[}person{]}.\par}
\medskip
\needspace{6\baselineskip}
\noindent\textbf{Mrf14}\enspace \citet{MerrifieldDanckert2014}.\par
{\small\noindent\textit{Design:} physiological / imaging.\quad \textit{Sample:} undergraduates.\quad \textit{Measure:} physiological measures during mood inductions.\par
\noindent\textit{Finding.} Induced boredom produced a physiological profile that differed from that of induced sadness.\par
\noindent\textit{Codes:} \texttt{distinct}.\quad \textit{Bearing:} (a) \textbf{N}, (b) \textbf{S}.\par
\noindent\textit{Rule trace:} (a) N: support {[}--{]}, challenge {[}--{]}; (b) S: motivation {[}--{]}, signal {[}distinct{]}, cost {[}--{]}.\par}
\medskip
\needspace{6\baselineskip}
\noindent\textbf{Wil14}\enspace \citet{WilsonEtAl2014}.\par
{\small\noindent\textit{Design:} experiment.\quad \textit{Sample:} undergraduates and community adults.\quad \textit{Measure:} enjoyment ratings; self-administered electric shocks.\par
\noindent\textit{Finding.} Participants generally did not enjoy spending several minutes alone with their thoughts, and many chose to give themselves mild electric shocks rather than sit without stimulation.\par
\noindent\textit{Codes:} \texttt{seek harm}.\quad \textit{Bearing:} (a) \textbf{N}, (b) \textbf{S/C}.\par
\noindent\textit{Rule trace:} (a) N: support {[}--{]}, challenge {[}--{]}; (b) S/C: motivation {[}seek{]}, signal {[}--{]}, cost {[}harm{]}.\par}
\medskip
\needspace{6\baselineskip}
\noindent\textbf{Hav15}\enspace \citet{Havermans2015}.\par
{\small\noindent\textit{Design:} experiment.\quad \textit{Sample:} laboratory sample.\quad \textit{Measure:} boredom induction; food intake; self-administered shocks.\par
\noindent\textit{Finding.} Under a boredom induction, participants ate more and administered more shocks to themselves than in a comparison condition.\par
\noindent\textit{Codes:} \texttt{seek harm}.\quad \textit{Bearing:} (a) \textbf{N}, (b) \textbf{S/C}.\par
\noindent\textit{Rule trace:} (a) N: support {[}--{]}, challenge {[}--{]}; (b) S/C: motivation {[}seek{]}, signal {[}--{]}, cost {[}harm{]}.\par}
\medskip
\needspace{6\baselineskip}
\noindent\textbf{Ned16}\enspace \citet{Nederkoorn2016}.\par
{\small\noindent\textit{Design:} experiment.\quad \textit{Sample:} laboratory sample.\quad \textit{Measure:} boring versus comparison film; self-administered shocks.\par
\noindent\textit{Finding.} A boring film led to more self-administered painful stimulation than a comparison film.\par
\noindent\textit{Codes:} \texttt{seek harm}.\quad \textit{Bearing:} (a) \textbf{N}, (b) \textbf{S/C}.\par
\noindent\textit{Rule trace:} (a) N: support {[}--{]}, challenge {[}--{]}; (b) S/C: motivation {[}seek{]}, signal {[}--{]}, cost {[}harm{]}.\par}
\medskip
\needspace{6\baselineskip}
\noindent\textbf{Str16}\enspace \citet{Struk2016}.\par
{\small\noindent\textit{Design:} correlational.\quad \textit{Sample:} undergraduates.\quad \textit{Measure:} BPS; regulatory-mode scales.\par
\noindent\textit{Finding.} Boredom proneness was associated with a regulatory profile of low locomotion and high assessment.\par
\noindent\textit{Codes:} \texttt{person}.\quad \textit{Bearing:} (a) \textbf{N}, (b) \textbf{N}.\par
\noindent\textit{Rule trace:} (a) N: support {[}--{]}, challenge {[}--{]}; (b) N: motivation {[}--{]}, signal {[}--{]}, cost {[}--{]}; accommodated by Cor.~3.4 of EXI-10-a: {[}person{]}.\par}
\medskip
\needspace{6\baselineskip}
\noindent\textbf{Til16}\enspace \citet{vanTilburgIgou2016}.\par
{\small\noindent\textit{Design:} experiment.\quad \textit{Sample:} student and adult samples.\quad \textit{Measure:} boredom inductions; political orientation.\par
\noindent\textit{Finding.} Boredom was associated with more extreme political orientations, interpreted as a way of re-establishing meaning.\par
\noindent\textit{Codes:} \texttt{meaningreg}.\quad \textit{Bearing:} (a) \textbf{N}, (b) \textbf{S}.\par
\noindent\textit{Rule trace:} (a) N: support {[}--{]}, challenge {[}--{]}; (b) S: motivation {[}meaningreg{]}, signal {[}--{]}, cost {[}--{]}.\par}
\medskip
\needspace{6\baselineskip}
\noindent\textbf{Tze16}\enspace \citet{Tze2016}.\par
{\small\noindent\textit{Design:} meta-analysis.\quad \textit{Sample:} school and university students (pooled).\quad \textit{Measure:} boredom and academic outcomes.\par
\noindent\textit{Finding.} Across studies, boredom was negatively related to academic outcomes.\par
\noindent\textit{Codes:} \texttt{perf-}.\quad \textit{Bearing:} (a) \textbf{N}, (b) \textbf{N}.\par
\noindent\textit{Rule trace:} (a) N: support {[}--{]}, challenge {[}--{]}; (b) N: motivation {[}--{]}, signal {[}--{]}, cost {[}--{]}; accommodated by Cor.~3.4 of EXI-10-a: {[}perf-{]}.\par}
\medskip
\needspace{6\baselineskip}
\noindent\textbf{Vod16}\enspace \citet{VodanovichWatt2016}.\par
{\small\noindent\textit{Design:} review.\quad \textit{Sample:} none (review).\quad \textit{Measure:} review of self-report boredom measures.\par
\noindent\textit{Finding.} Updates the review of trait and state boredom scales, including newer instruments.\par
\noindent\textit{Codes:} \texttt{meas}.\quad \textit{Bearing:} (a) \textbf{N}, (b) \textbf{N}.\par
\noindent\textit{Rule trace:} (a) N: support {[}--{]}, challenge {[}--{]}; (b) N: motivation {[}--{]}, signal {[}--{]}, cost {[}--{]}.\par}
\medskip
\needspace{6\baselineskip}
\noindent\textbf{Chi17}\enspace \citet{Chin2017}.\par
{\small\noindent\textit{Design:} experience sampling.\quad \textit{Sample:} adults (smartphone sampling).\quad \textit{Measure:} momentary boredom reports.\par
\noindent\textit{Finding.} Boredom was common in everyday life and varied with activity and social context.\par
\noindent\textit{Codes:} \texttt{relational}.\quad \textit{Bearing:} (a) \textbf{S}, (b) \textbf{N}.\par
\noindent\textit{Rule trace:} (a) S: support {[}relational{]}, challenge {[}--{]}; (b) N: motivation {[}--{]}, signal {[}--{]}, cost {[}--{]}.\par}
\medskip
\needspace{6\baselineskip}
\noindent\textbf{Isa17}\enspace \citet{Isacescu2017}.\par
{\small\noindent\textit{Design:} correlational.\quad \textit{Sample:} undergraduates.\quad \textit{Measure:} boredom proneness; self-control and affect measures.\par
\noindent\textit{Finding.} Self-control and affective variables predicted boredom proneness.\par
\noindent\textit{Codes:} \texttt{person}.\quad \textit{Bearing:} (a) \textbf{N}, (b) \textbf{N}.\par
\noindent\textit{Rule trace:} (a) N: support {[}--{]}, challenge {[}--{]}; (b) N: motivation {[}--{]}, signal {[}--{]}, cost {[}--{]}; accommodated by Cor.~3.4 of EXI-10-a: {[}person{]}.\par}
\medskip
\needspace{6\baselineskip}
\noindent\textbf{Str17}\enspace \citet{Struk2017}.\par
{\small\noindent\textit{Design:} psychometric.\quad \textit{Sample:} undergraduates and online adults.\quad \textit{Measure:} Short Boredom Proneness Scale (SBPS).\par
\noindent\textit{Finding.} Develops an eight-item unidimensional short form of the boredom proneness construct with good psychometric properties.\par
\noindent\textit{Codes:} \texttt{meas}.\quad \textit{Bearing:} (a) \textbf{N}, (b) \textbf{N}.\par
\noindent\textit{Rule trace:} (a) N: support {[}--{]}, challenge {[}--{]}; (b) N: motivation {[}--{]}, signal {[}--{]}, cost {[}--{]}.\par}
\medskip
\needspace{6\baselineskip}
\noindent\textbf{Til17}\enspace \citet{vanTilburgIgou2017}.\par
{\small\noindent\textit{Design:} experiment.\quad \textit{Sample:} student and adult samples.\quad \textit{Measure:} recalled emotions; appraisal and motivation ratings.\par
\noindent\textit{Finding.} Boredom differed from other negative emotions in its profile, notably in perceived meaninglessness and in the wish to pursue a different goal.\par
\noindent\textit{Codes:} \texttt{meaning+ signal switch}.\quad \textit{Bearing:} (a) \textbf{S}, (b) \textbf{S}.\par
\noindent\textit{Rule trace:} (a) S: support {[}meaning+{]}, challenge {[}--{]}; (b) S: motivation {[}switch{]}, signal {[}signal{]}, cost {[}--{]}.\par}
\medskip
\needspace{6\baselineskip}
\noindent\textbf{DnH18}\enspace \citet{Danckert2018restless}.\par
{\small\noindent\textit{Design:} correlational.\quad \textit{Sample:} undergraduates and online adults.\quad \textit{Measure:} state boredom; arousal and restlessness reports.\par
\noindent\textit{Finding.} Boredom was associated with low arousal combined with restlessness.\par
\noindent\textit{Codes:} \texttt{distinct}.\quad \textit{Bearing:} (a) \textbf{N}, (b) \textbf{S}.\par
\noindent\textit{Rule trace:} (a) N: support {[}--{]}, challenge {[}--{]}; (b) S: motivation {[}--{]}, signal {[}distinct{]}, cost {[}--{]}.\par}
\medskip
\needspace{6\baselineskip}
\noindent\textbf{DnM18}\enspace \citet{DanckertMerrifield2018}.\par
{\small\noindent\textit{Design:} physiological / imaging.\quad \textit{Sample:} adult participants.\quad \textit{Measure:} fMRI during boredom induction.\par
\noindent\textit{Finding.} Boredom induction was associated with activity in the default mode network.\par
\noindent\textit{Codes:} \texttt{neuro}.\quad \textit{Bearing:} (a) \textbf{N}, (b) \textbf{N}.\par
\noindent\textit{Rule trace:} (a) N: support {[}--{]}, challenge {[}--{]}; (b) N: motivation {[}--{]}, signal {[}--{]}, cost {[}--{]}.\par}
\medskip
\needspace{6\baselineskip}
\noindent\textbf{Dnk18}\enspace \citet{DanckertMugon2018}.\par
{\small\noindent\textit{Design:} theory.\quad \textit{Sample:} none (theory).\quad \textit{Measure:} --.\par
\noindent\textit{Finding.} Argues that boredom signals that one's current engagement is unsatisfactory and prompts the exploration of alternatives.\par
\noindent\textit{Codes:} \texttt{signal seek}.\quad \textit{Bearing:} (a) \textbf{N}, (b) \textbf{S}.\par
\noindent\textit{Rule trace:} (a) N: support {[}--{]}, challenge {[}--{]}; (b) S: motivation {[}seek{]}, signal {[}signal{]}, cost {[}--{]}.\par}
\medskip
\needspace{6\baselineskip}
\noindent\textbf{Elp18}\enspace \citet{Elpidorou2018}.\par
{\small\noindent\textit{Design:} theory.\quad \textit{Sample:} none (theory).\quad \textit{Measure:} --.\par
\noindent\textit{Finding.} Argues that boredom is valuable as a signal of a mismatch between the agent's situation and her desires, and as a motivation to change it.\par
\noindent\textit{Codes:} \texttt{signal switch}.\quad \textit{Bearing:} (a) \textbf{N}, (b) \textbf{S}.\par
\noindent\textit{Rule trace:} (a) N: support {[}--{]}, challenge {[}--{]}; (b) S: motivation {[}switch{]}, signal {[}signal{]}, cost {[}--{]}.\par}
\medskip
\needspace{6\baselineskip}
\noindent\textbf{Hun18}\enspace \citet{HunterEastwood2018}.\par
{\small\noindent\textit{Design:} experiment.\quad \textit{Sample:} undergraduates.\quad \textit{Measure:} trait and state boredom; sustained-attention task.\par
\noindent\textit{Finding.} Examines the relations of trait and state boredom to sustained attention; attentional performance and boredom were related.\par
\noindent\textit{Codes:} \texttt{attn}.\quad \textit{Bearing:} (a) \textbf{S}, (b) \textbf{N}.\par
\noindent\textit{Rule trace:} (a) S: support {[}attn{]}, challenge {[}--{]}; (b) N: motivation {[}--{]}, signal {[}--{]}, cost {[}--{]}.\par}
\medskip
\needspace{6\baselineskip}
\noindent\textbf{Mug18}\enspace \citet{Mugon2018}.\par
{\small\noindent\textit{Design:} correlational.\quad \textit{Sample:} undergraduates.\quad \textit{Measure:} BPS; regulatory modes.\par
\noindent\textit{Finding.} Boredom proneness was associated with difficulty in initiating action (low locomotion).\par
\noindent\textit{Codes:} \texttt{person}.\quad \textit{Bearing:} (a) \textbf{N}, (b) \textbf{N}.\par
\noindent\textit{Rule trace:} (a) N: support {[}--{]}, challenge {[}--{]}; (b) N: motivation {[}--{]}, signal {[}--{]}, cost {[}--{]}; accommodated by Cor.~3.4 of EXI-10-a: {[}person{]}.\par}
\medskip
\needspace{6\baselineskip}
\noindent\textbf{Raf18}\enspace \citet{Raffaelli2018}.\par
{\small\noindent\textit{Design:} review.\quad \textit{Sample:} none (review).\quad \textit{Measure:} --.\par
\noindent\textit{Finding.} Reviews what is and is not known about boredom across definitional, measurement and neuroscientific work.\par
\noindent\textit{Codes:} \texttt{meas}.\quad \textit{Bearing:} (a) \textbf{N}, (b) \textbf{N}.\par
\noindent\textit{Rule trace:} (a) N: support {[}--{]}, challenge {[}--{]}; (b) N: motivation {[}--{]}, signal {[}--{]}, cost {[}--{]}.\par}
\medskip
\needspace{6\baselineskip}
\noindent\textbf{Wes18}\enspace \citet{WestgateWilson2018}.\par
{\small\noindent\textit{Design:} theory.\quad \textit{Sample:} none (theory).\quad \textit{Measure:} --.\par
\noindent\textit{Finding.} The MAC model: boredom arises from a mismatch between cognitive demands and available resources (too low or too high) and from a lack of meaning, both signaling attentional or value deficits.\par
\noindent\textit{Codes:} \texttt{under+ over+ meaning+ attn signal}.\quad \textit{Bearing:} (a) \textbf{S}, (b) \textbf{S}.\par
\noindent\textit{Rule trace:} (a) S: support {[}attn, meaning+, over+, under+{]}, challenge {[}--{]}; (b) S: motivation {[}--{]}, signal {[}signal{]}, cost {[}--{]}.\par}
\medskip
\needspace{6\baselineskip}
\noindent\textbf{Ben19}\enspace \citet{BenchLench2019}.\par
{\small\noindent\textit{Design:} experiment.\quad \textit{Sample:} undergraduates.\quad \textit{Measure:} boredom induction; choice of novel stimuli.\par
\noindent\textit{Finding.} Induced boredom increased the choice of novel experiences, even unpleasant ones.\par
\noindent\textit{Codes:} \texttt{seek}.\quad \textit{Bearing:} (a) \textbf{N}, (b) \textbf{S}.\par
\noindent\textit{Rule trace:} (a) N: support {[}--{]}, challenge {[}--{]}; (b) S: motivation {[}seek{]}, signal {[}--{]}, cost {[}--{]}.\par}
\medskip
\needspace{6\baselineskip}
\noindent\textbf{DnE20}\enspace \citet{DanckertEastwood2020}.\par
{\small\noindent\textit{Design:} review.\quad \textit{Sample:} none (review).\quad \textit{Measure:} --.\par
\noindent\textit{Finding.} Synthesizes research on boredom as a call to action: a signal that one is not engaged that prompts the search for engagement.\par
\noindent\textit{Codes:} \texttt{attn signal seek}.\quad \textit{Bearing:} (a) \textbf{S}, (b) \textbf{S}.\par
\noindent\textit{Rule trace:} (a) S: support {[}attn{]}, challenge {[}--{]}; (b) S: motivation {[}seek{]}, signal {[}signal{]}, cost {[}--{]}.\par}
\medskip
\needspace{6\baselineskip}
\noindent\textbf{Wol20}\enspace \citet{Wolff2020}.\par
{\small\noindent\textit{Design:} correlational.\quad \textit{Sample:} adults (online).\quad \textit{Measure:} boredom proneness; self-control; adherence to distancing rules.\par
\noindent\textit{Finding.} Higher boredom proneness and lower self-control were associated with weaker adherence to social-distancing guidelines.\par
\noindent\textit{Codes:} \texttt{chronic}.\quad \textit{Bearing:} (a) \textbf{N}, (b) \textbf{N}.\par
\noindent\textit{Rule trace:} (a) N: support {[}--{]}, challenge {[}--{]}; (b) N: motivation {[}--{]}, signal {[}--{]}, cost {[}--{]}; accommodated by Cor.~3.4 of EXI-10-a: {[}chronic{]}.\par}
\medskip
\needspace{6\baselineskip}
\noindent\textbf{Boy21}\enspace \citet{Boylan2021}.\par
{\small\noindent\textit{Design:} correlational.\quad \textit{Sample:} adults (online).\quad \textit{Measure:} boredom proneness; desire to act; rule-breaking.\par
\noindent\textit{Finding.} Trait boredom proneness was associated with a desire to act and with self-reported breaking of pandemic rules.\par
\noindent\textit{Codes:} \texttt{chronic seek}.\quad \textit{Bearing:} (a) \textbf{N}, (b) \textbf{S}.\par
\noindent\textit{Rule trace:} (a) N: support {[}--{]}, challenge {[}--{]}; (b) S: motivation {[}seek{]}, signal {[}--{]}, cost {[}--{]}; accommodated by Cor.~3.4 of EXI-10-a: {[}chronic{]}.\par}
\medskip
\needspace{6\baselineskip}
\noindent\textbf{Tam21}\enspace \citet{Tam2021}.\par
{\small\noindent\textit{Design:} theory.\quad \textit{Sample:} none (theory).\quad \textit{Measure:} --.\par
\noindent\textit{Finding.} The Boredom Feedback Model: boredom is feedback that attention is not adequately engaged, prompting attempts to re-engage.\par
\noindent\textit{Codes:} \texttt{attn signal seek}.\quad \textit{Bearing:} (a) \textbf{S}, (b) \textbf{S}.\par
\noindent\textit{Rule trace:} (a) S: support {[}attn{]}, challenge {[}--{]}; (b) S: motivation {[}seek{]}, signal {[}signal{]}, cost {[}--{]}.\par}
\medskip


\section{Computed tables}\label{app:b}

\paragraph{Codes.} Table~\ref{tab:codes} lists the codes with the number of entries carrying each.

\begin{table}[h]\caption{Codes and their frequencies.}\label{tab:codes}
\begin{center}\small
\begin{tabular}{@{}llr@{}}\toprule
Code & Meaning & Studies \\ \midrule
\texttt{under+} & boredom rises with under-challenge or monotony & 7 \\
\texttt{over+} & boredom rises with over-challenge & 1 \\
\texttt{meaning+} & boredom rises with lack of meaning or value & 5 \\
\texttt{fit-} & matched challenge and skill yields engagement & 6 \\
\texttt{attn} & boredom linked to failure of attentional engagement & 7 \\
\texttt{relational} & state boredom depends on person and situation jointly & 2 \\
\texttt{overanx} & over-challenge yields anxiety, not boredom & 2 \\
\texttt{seek} & boredom prompts seeking of stimulation or novelty & 11 \\
\texttt{switch} & boredom prompts change of goal or of manner of engagement & 6 \\
\texttt{meaningreg} & boredom prompts meaning-restoring responses & 3 \\
\texttt{creat} & boredom followed by more divergent thinking & 1 \\
\texttt{signal} & boredom described as a signal or feedback state & 9 \\
\texttt{distinct} & boredom has a distinct experiential or physiological profile & 4 \\
\texttt{harm} & boredom prompts self-harming or costly responses & 3 \\
\texttt{perf-} & boredom related to lower performance on the task & 2 \\
\texttt{chronic} & trait proneness related to poor outcomes & 4 \\
\texttt{person} & trait or individual-difference correlates & 9 \\
\texttt{meas} & measurement or conceptual work & 9 \\
\texttt{neuro} & neural correlates & 1 \\
\bottomrule\end{tabular}\end{center}

\end{table}

\paragraph{Overview.} One line per entry, in order of publication, with the computed bearings under both rules.

{\footnotesize\setlength{\tabcolsep}{3pt}
\begin{longtable}{@{}lrllccc@{}}
\toprule
Id & Year & Design & Codes & (a) & (b) & (b) strict \\
\midrule\endhead
\bottomrule\endfoot
Bex54 & 1954 & experiment & \texttt{under+ seek} & S & S & S \\
Ber60 & 1960 & theory & \texttt{under+ seek} & S & S & S \\
Csi75 & 1975 & interview-based & \texttt{under+ fit- overanx} & S/C & N & N \\
Zuc79 & 1979 & psychometric & \texttt{meas person} & N & N & N \\
OHa81 & 1981 & review & \texttt{under+ attn} & S & N & N \\
Far86 & 1986 & psychometric & \texttt{meas chronic} & N & N & C \\
Csi89 & 1989 & experience sampling & \texttt{fit-} & S & N & N \\
Csi90 & 1990 & theory & \texttt{under+ fit- overanx} & S/C & N & N \\
Vod90 & 1990 & psychometric & \texttt{meas} & N & N & N \\
Lar91 & 1991 & experience sampling & \texttt{person} & N & N & N \\
San92 & 1992 & experiment & \texttt{switch} & N & S & S \\
Mon96 & 1996 & experience sampling & \texttt{fit-} & S & N & N \\
Som00 & 2000 & correlational & \texttt{chronic} & N & N & C \\
Vod03 & 2003 & review & \texttt{meas} & N & N & N \\
Dan05 & 2005 & correlational & \texttt{distinct} & N & S & S \\
Eas07 & 2007 & correlational & \texttt{person} & N & N & N \\
Eng08 & 2008 & correlational & \texttt{fit-} & S & N & N \\
Fah09 & 2009 & correlational & \texttt{meaning+} & S & N & N \\
Pek10 & 2010 & correlational & \texttt{meaning+ perf-} & S & N & C \\
Gol11 & 2011 & psychometric & \texttt{distinct} & N & S & S \\
Kel11 & 2011 & experiment & \texttt{fit-} & S & N & N \\
Til11 & 2011 & experiment & \texttt{meaningreg} & N & S & S \\
Eas12 & 2012 & theory & \texttt{attn signal} & S & S & S \\
Mal12 & 2012 & experiment & \texttt{attn person} & S & N & N \\
Til12 & 2012 & experiment & \texttt{under+ meaning+ signal switch} & S & S & S \\
Ben13 & 2013 & theory & \texttt{switch} & N & S & S \\
Fah13 & 2013 & psychometric & \texttt{meas} & N & N & N \\
Til13 & 2013 & experiment & \texttt{meaningreg} & N & S & S \\
Elp14 & 2014 & theory & \texttt{signal switch} & N & S & S \\
Ger14 & 2014 & correlational & \texttt{person} & N & N & N \\
Goe14 & 2014 & experience sampling & \texttt{meas seek} & N & S & S \\
Man14 & 2014 & experiment & \texttt{creat} & N & S & S \\
Mer14 & 2014 & correlational & \texttt{relational person} & S & N & N \\
Mrf14 & 2014 & physiological / imaging & \texttt{distinct} & N & S & S \\
Wil14 & 2014 & experiment & \texttt{seek harm} & N & S/C & S/C \\
Hav15 & 2015 & experiment & \texttt{seek harm} & N & S/C & S/C \\
Ned16 & 2016 & experiment & \texttt{seek harm} & N & S/C & S/C \\
Str16 & 2016 & correlational & \texttt{person} & N & N & N \\
Til16 & 2016 & experiment & \texttt{meaningreg} & N & S & S \\
Tze16 & 2016 & meta-analysis & \texttt{perf-} & N & N & C \\
Vod16 & 2016 & review & \texttt{meas} & N & N & N \\
Chi17 & 2017 & experience sampling & \texttt{relational} & S & N & N \\
Isa17 & 2017 & correlational & \texttt{person} & N & N & N \\
Str17 & 2017 & psychometric & \texttt{meas} & N & N & N \\
Til17 & 2017 & experiment & \texttt{meaning+ signal switch} & S & S & S \\
DnH18 & 2018 & correlational & \texttt{distinct} & N & S & S \\
DnM18 & 2018 & physiological / imaging & \texttt{neuro} & N & N & N \\
Dnk18 & 2018 & theory & \texttt{signal seek} & N & S & S \\
Elp18 & 2018 & theory & \texttt{signal switch} & N & S & S \\
Hun18 & 2018 & experiment & \texttt{attn} & S & N & N \\
Mug18 & 2018 & correlational & \texttt{person} & N & N & N \\
Raf18 & 2018 & review & \texttt{meas} & N & N & N \\
Wes18 & 2018 & theory & \texttt{under+ over+ meaning+ attn signal} & S & S & S \\
Ben19 & 2019 & experiment & \texttt{seek} & N & S & S \\
DnE20 & 2020 & review & \texttt{attn signal seek} & S & S & S \\
Wol20 & 2020 & correlational & \texttt{chronic} & N & N & C \\
Boy21 & 2021 & correlational & \texttt{chronic seek} & N & S & S/C \\
Tam21 & 2021 & theory & \texttt{attn signal seek} & S & S & S \\
\end{longtable}}


\paragraph{Bearing on (a) by bearing on (b).} All entries, and empirical entries only.

\begin{table}[h]\caption{Cross-tabulation of the bearings, all entries.}\label{tab:ab}
\begin{center}\small
\begin{tabular}{@{}lrrrrr@{}}\toprule
(a) $\backslash$ (b) & S & S/C & N & C & Total \\ \midrule
S & 8 & . & 11 & . & 19 \\
S/C & . & . & 2 & . & 2 \\
N & 16 & 3 & 18 & . & 37 \\
C & . & . & . & . & 0 \\
\midrule Total & 24 & 3 & 31 & 0 & 58 \\
\bottomrule\end{tabular}\end{center}

\end{table}

\begin{table}[h]\caption{Cross-tabulation of the bearings, empirical entries only.}\label{tab:abemp}
\begin{center}\small
\begin{tabular}{@{}lrrrrr@{}}\toprule
(a) $\backslash$ (b) & S & S/C & N & C & Total \\ \midrule
S & 3 & . & 10 & . & 13 \\
S/C & . & . & 1 & . & 1 \\
N & 12 & 3 & 15 & . & 30 \\
C & . & . & . & . & 0 \\
\midrule Total & 15 & 3 & 26 & 0 & 44 \\
\bottomrule\end{tabular}\end{center}

\end{table}

\paragraph{By design and by period.}

\begin{table}[h]\caption{Bearing on the mismatch thesis by design.}\label{tab:da}
\begin{center}\small
\begin{tabular}{@{}lrrrrr@{}}\toprule
Design (a) & S & S/C & N & C & Total \\ \midrule
experiment & 6 & . & 9 & . & 15 \\
experience sampling & 3 & . & 2 & . & 5 \\
correlational & 4 & . & 10 & . & 14 \\
psychometric & . & . & 6 & . & 6 \\
physiological / imaging & . & . & 2 & . & 2 \\
meta-analysis & . & . & 1 & . & 1 \\
review & 2 & . & 3 & . & 5 \\
theory & 4 & 1 & 4 & . & 9 \\
interview-based & . & 1 & . & . & 1 \\
\midrule Total & 19 & 2 & 37 & 0 & 58 \\
\bottomrule\end{tabular}\end{center}

\end{table}

\begin{table}[h]\caption{Bearing on the function thesis by design (lenient rule).}\label{tab:db}
\begin{center}\small
\begin{tabular}{@{}lrrrrr@{}}\toprule
Design (b) & S & S/C & N & C & Total \\ \midrule
experiment & 9 & 3 & 3 & . & 15 \\
experience sampling & 1 & . & 4 & . & 5 \\
correlational & 3 & . & 11 & . & 14 \\
psychometric & 1 & . & 5 & . & 6 \\
physiological / imaging & 1 & . & 1 & . & 2 \\
meta-analysis & . & . & 1 & . & 1 \\
review & 1 & . & 4 & . & 5 \\
theory & 8 & . & 1 & . & 9 \\
interview-based & . & . & 1 & . & 1 \\
\midrule Total & 24 & 3 & 31 & 0 & 58 \\
\bottomrule\end{tabular}\end{center}

\end{table}

\begin{table}[h]\caption{Bearings by period of publication.}\label{tab:period}
\begin{center}\small
\begin{tabular}{@{}lrrrrrrrrr@{}}\toprule
 & \multicolumn{4}{c}{(a)} & \multicolumn{4}{c}{(b)} & \\
Period & S & S/C & N & C & S & S/C & N & C & Total \\ \midrule
before 2000 & 5 & 2 & 5 & . & 3 & . & 9 & . & 12 \\
2000s & 2 & . & 4 & . & 1 & . & 5 & . & 6 \\
2010s & 10 & . & 26 & . & 17 & 3 & 16 & . & 36 \\
2020s & 2 & . & 2 & . & 3 & . & 1 & . & 4 \\
\bottomrule\end{tabular}\end{center}

\end{table}

\paragraph{Sensitivity to the rule for costs.} Table~\ref{tab:strict} cross-tabulates the bearing on (b) under the lenient and the strict rule. Entries move only from S to S/C or from N to C; the motivational evidence is unaffected.

\begin{table}[h]\caption{Function thesis: lenient against strict rule.}\label{tab:strict}
\begin{center}\small
\begin{tabular}{@{}lrrrrr@{}}\toprule
(b) lenient $\backslash$ strict & S & S/C & N & C & Total \\ \midrule
S & 23 & 1 & . & . & 24 \\
S/C & . & 3 & . & . & 3 \\
N & . & . & 26 & 5 & 31 \\
C & . & . & . & . & 0 \\
\midrule Total & 23 & 4 & 26 & 5 & 58 \\
\bottomrule\end{tabular}\end{center}

\end{table}

\paragraph{Incidence of statements.} Table~\ref{tab:matrix} gives, for each entry, which statements of EXI-10-a its codes bear on: $+$ for, $-$ against, $\pm$ both, with the code-to-statement links of Table~\ref{tab:clauses}. The last row counts the entries bearing on each statement.

{\footnotesize\setlength{\tabcolsep}{4pt}
\begin{longtable}{@{}lcccccccc@{}}
\caption{Entries and the statements of EXI-10-a they bear on.}\label{tab:matrix}\\
\toprule
Id & Und & Ovr & Mng & A2.4 & L3.1 & L3.2 & T3.3 & C3.4 \\
\midrule\endhead
\bottomrule\endfoot
Bex54 & $+$ & . & . & . & . & $+$ & . & . \\
Ber60 & $+$ & . & . & . & . & $+$ & . & . \\
Csi75 & $+$ & $-$ & . & $+$ & . & . & $+$ & . \\
Zuc79 & . & . & . & . & . & . & . & $+$ \\
OHa81 & $+$ & . & . & . & $+$ & . & . & . \\
Far86 & . & . & . & . & . & . & . & $+$ \\
Csi89 & . & . & . & $+$ & . & . & $+$ & . \\
Csi90 & $+$ & $-$ & . & $+$ & . & . & $+$ & . \\
Vod90 & . & . & . & . & . & . & . & . \\
Lar91 & . & . & . & . & . & . & . & $+$ \\
San92 & . & . & . & . & . & $+$ & $+$ & . \\
Mon96 & . & . & . & $+$ & . & . & $+$ & . \\
Som00 & . & . & . & . & . & . & . & $+$ \\
Vod03 & . & . & . & . & . & . & . & . \\
Dan05 & . & . & . & . & $+$ & . & . & . \\
Eas07 & . & . & . & . & . & . & . & $+$ \\
Eng08 & . & . & . & $+$ & . & . & $+$ & . \\
Fah09 & . & . & $+$ & . & . & . & . & . \\
Pek10 & . & . & $+$ & . & . & . & . & . \\
Gol11 & . & . & . & . & $+$ & . & . & . \\
Kel11 & . & . & . & $+$ & . & . & $+$ & . \\
Til11 & . & . & . & . & . & $+$ & $+$ & . \\
Eas12 & . & . & . & . & $+$ & . & . & . \\
Mal12 & . & . & . & . & $+$ & . & . & $+$ \\
Til12 & $+$ & . & $+$ & . & $+$ & $+$ & $+$ & . \\
Ben13 & . & . & . & . & . & $+$ & $+$ & . \\
Fah13 & . & . & . & . & . & . & . & . \\
Til13 & . & . & . & . & . & $+$ & $+$ & . \\
Elp14 & . & . & . & . & $+$ & $+$ & $+$ & . \\
Ger14 & . & . & . & . & . & . & . & $+$ \\
Goe14 & . & . & . & . & . & $+$ & . & . \\
Man14 & . & . & . & . & . & $+$ & $+$ & . \\
Mer14 & . & . & . & . & . & . & . & $+$ \\
Mrf14 & . & . & . & . & $+$ & . & . & . \\
Wil14 & . & . & . & . & . & $+$ & $-$ & . \\
Hav15 & . & . & . & . & . & $+$ & $-$ & . \\
Ned16 & . & . & . & . & . & $+$ & $-$ & . \\
Str16 & . & . & . & . & . & . & . & $+$ \\
Til16 & . & . & . & . & . & $+$ & $+$ & . \\
Tze16 & . & . & . & . & . & . & . & . \\
Vod16 & . & . & . & . & . & . & . & . \\
Chi17 & . & . & . & . & . & . & . & . \\
Isa17 & . & . & . & . & . & . & . & $+$ \\
Str17 & . & . & . & . & . & . & . & . \\
Til17 & . & . & $+$ & . & $+$ & $+$ & $+$ & . \\
DnH18 & . & . & . & . & $+$ & . & . & . \\
DnM18 & . & . & . & . & . & . & . & . \\
Dnk18 & . & . & . & . & $+$ & $+$ & . & . \\
Elp18 & . & . & . & . & $+$ & $+$ & $+$ & . \\
Hun18 & . & . & . & . & $+$ & . & . & . \\
Mug18 & . & . & . & . & . & . & . & $+$ \\
Raf18 & . & . & . & . & . & . & . & . \\
Wes18 & $+$ & $+$ & $+$ & . & $+$ & . & . & . \\
Ben19 & . & . & . & . & . & $+$ & . & . \\
DnE20 & . & . & . & . & $+$ & $+$ & . & . \\
Wol20 & . & . & . & . & . & . & . & $+$ \\
Boy21 & . & . & . & . & . & $+$ & . & $+$ \\
Tam21 & . & . & . & . & $+$ & $+$ & . & . \\
\midrule Entries & 7 & 3 & 5 & 6 & 16 & 21 & 19 & 13 \\
\end{longtable}}


\paragraph{Samples.} Table~\ref{tab:samples} gives the bearings by kind of sample, the samples being grouped by the program from the descriptions in the entries.

\begin{table}[h]\caption{Bearings by kind of sample.}\label{tab:samples}
\begin{center}\small
\begin{tabular}{@{}lrrrrrrrrr@{}}\toprule
 & \multicolumn{4}{c}{(a)} & \multicolumn{4}{c}{(b)} & \\
Sample & S & S/C & N & C & S & S/C & N & C & Total \\ \midrule
students & 9 & . & 21 & . & 12 & 1 & 17 & . & 30 \\
school students & 1 & . & 1 & . & . & . & 2 & . & 2 \\
online or smartphone adults & 1 & . & 2 & . & 1 & . & 2 & . & 3 \\
other adult or laboratory & 2 & 1 & 6 & . & 2 & 2 & 5 & . & 9 \\
none (theory/review) & 6 & 1 & 7 & . & 9 & . & 5 & . & 14 \\
\bottomrule\end{tabular}\end{center}

\end{table}

\paragraph{Statements of EXI-10-a.} Table~\ref{tab:clauses} lists, for each clause of Definition~2.2 and each of the main statements of EXI-10-a, the entries whose codes bear on it, with their number and the number of empirical ones. A statement is linked to an entry by the following codes: under-challenge \texttt{under+}; over-challenge \texttt{over+} (for) and \texttt{overanx} (against); meaninglessness \texttt{meaning+}; Assumption~2.4 \texttt{fit-}; Lemma~3.1 \texttt{signal}, \texttt{distinct}, \texttt{attn}; Lemma~3.2 the motivational codes; Theorem~3.3 \texttt{fit-} and the codes of goal change, meaning regulation and creativity (for) and \texttt{harm} (against); Corollary~3.4 \texttt{chronic} and \texttt{person}.

\begin{center}\refstepcounter{table}\label{tab:clauses}Table~\thetable: Statements of EXI-10-a and the entries bearing on them.\end{center}
{\small\setlength{\tabcolsep}{3pt}
\begin{longtable}{@{}p{0.27\linewidth}p{0.43\linewidth}p{0.24\linewidth}@{}}
\toprule
Statement of EXI-10-a & Bearing for (all study ids) & Bearing against \\
\midrule\endhead
\bottomrule\endfoot
Def.~2.2, under-challenge & Bex54, Ber60, Csi75, OHa81, Csi90, Til12, Wes18 (7; 3 empirical) & -- \\
Def.~2.2, over-challenge & Wes18 (1; 0 empirical) & Csi75, Csi90 \\
Def.~2.2, meaninglessness & Fah09, Pek10, Til12, Til17, Wes18 (5; 4 empirical) & -- \\
Assumption~2.4 (value of engagement) & Csi75, Csi89, Csi90, Mon96, Eng08, Kel11 (6; 5 empirical) & -- \\
Lemma~3.1 (signal) & OHa81, Dan05, Gol11, Eas12, Mal12, Til12, Elp14, Mrf14, Til17, DnH18, Dnk18, Elp18, Hun18, Wes18, DnE20, Tam21 (16; 8 empirical) & -- \\
Lemma~3.2 (motivation) & Bex54, Ber60, San92, Til11, Til12, Ben13, Til13, Elp14, Goe14, Man14, Wil14, Hav15, Ned16, Til16, Til17, Dnk18, Elp18, Ben19, DnE20, Boy21, Tam21 (21; 14 empirical) & -- \\
Theorem~3.3 (adaptive function) & Csi75, Csi89, Csi90, San92, Mon96, Eng08, Kel11, Til11, Til12, Ben13, Til13, Elp14, Man14, Til16, Til17, Elp18 (16; 12 empirical) & Wil14, Hav15, Ned16 \\
Corollary~3.4 (chronic boredom) & Zuc79, Far86, Lar91, Som00, Eas07, Mal12, Ger14, Mer14, Str16, Isa17, Mug18, Wol20, Boy21 (13; 13 empirical) & -- \\
\end{longtable}}


\begin{thebibliography}{99}\small
\bibitem[Bench and Lench(2013)]{BenchLench2013} Bench, S.~W. and Lench, H.~C. (2013). On the function of boredom. \emph{Behavioral Sciences}, 3(3), 459--472.
\bibitem[Bench and Lench(2019)]{BenchLench2019} Bench, S.~W. and Lench, H.~C. (2019). Boredom as a seeking state: Boredom prompts the pursuit of novel (even negative) experiences. \emph{Emotion}, 19(2), 242--254.
\bibitem[Berlyne(1960)]{Berlyne1960} Berlyne, D.~E. (1960). \emph{Conflict, Arousal, and Curiosity}. McGraw-Hill.
\bibitem[Bexton et~al.(1954)]{Bexton1954} Bexton, W.~H., Heron, W. and Scott, T.~H. (1954). Effects of decreased variation in the sensory environment. \emph{Canadian Journal of Psychology}, 8(2), 70--76.
\bibitem[Boylan et~al.(2021)]{Boylan2021} Boylan, J., Seli, P., Scholer, A.~A. and Danckert, J. (2021). Boredom in the COVID-19 pandemic: Trait boredom proneness, the desire to act, and rule-breaking. \emph{Personality and Individual Differences}, 171, 110387.
\bibitem[Chin et~al.(2017)]{Chin2017} Chin, A., Markey, A., Bhargava, S., Kassam, K.~S. and Loewenstein, G. (2017). Bored in the USA: Experience sampling and boredom in everyday life. \emph{Emotion}, 17(2).
\bibitem[Csikszentmihalyi(1975)]{Csikszentmihalyi1975} Csikszentmihalyi, M. (1975). \emph{Beyond Boredom and Anxiety}. Jossey-Bass.
\bibitem[Csikszentmihalyi(1990)]{Csikszentmihalyi1990} Csikszentmihalyi, M. (1990). \emph{Flow: The Psychology of Optimal Experience}. Harper \& Row.
\bibitem[Csikszentmihalyi and LeFevre(1989)]{CsikszentmihalyiLeFevre1989} Csikszentmihalyi, M. and LeFevre, J. (1989). Optimal experience in work and leisure. \emph{Journal of Personality and Social Psychology}, 56(5).
\bibitem[Danckert and Allman(2005)]{DanckertAllman2005} Danckert, J.~A. and Allman, A.-A.~A. (2005). Time flies when you're having fun: Temporal estimation and the experience of boredom. \emph{Brain and Cognition}, 59(3).
\bibitem[Danckert and Eastwood(2020)]{DanckertEastwood2020} Danckert, J. and Eastwood, J.~D. (2020). \emph{Out of My Skull: The Psychology of Boredom}. Harvard University Press.
\bibitem[Danckert et~al.(2018a)]{Danckert2018restless} Danckert, J., Hammerschmidt, T., Marty-Dugas, J. and Smilek, D. (2018a). Boredom: Under-aroused and restless. \emph{Consciousness and Cognition}, 61.
\bibitem[Danckert and Merrifield(2018)]{DanckertMerrifield2018} Danckert, J. and Merrifield, C. (2018). Boredom, sustained attention and the default mode network. \emph{Experimental Brain Research}, 236(9).
\bibitem[Danckert et~al.(2018b)]{DanckertMugon2018} Danckert, J., Mugon, J., Struk, A. and Eastwood, J. (2018b). Boredom: What is it good for? In H.~C. Lench (Ed.), \emph{The Function of Emotions: When and Why Emotions Help Us}. Springer.
\bibitem[Eastwood et~al.(2007)]{Eastwood2007} Eastwood, J.~D., Cavaliere, C., Fahlman, S.~A. and Eastwood, A.~E. (2007). A desire for desires: Boredom and its relation to alexithymia. \emph{Personality and Individual Differences}, 42(6).
\bibitem[Eastwood et~al.(2012)]{EastwoodEtAl2012} Eastwood, J.~D., Frischen, A., Fenske, M.~J. and Smilek, D. (2012). The unengaged mind: Defining boredom in terms of attention. \emph{Perspectives on Psychological Science}, 7(5), 482--495.
\bibitem[Elpidorou(2014)]{Elpidorou2014} Elpidorou, A. (2014). The bright side of boredom. \emph{Frontiers in Psychology}, 5, 1245.
\bibitem[Elpidorou(2018)]{Elpidorou2018} Elpidorou, A. (2018). The good of boredom. \emph{Philosophical Psychology}, 31(3), 323--351.
\bibitem[Engeser and Rheinberg(2008)]{EngeserRheinberg2008} Engeser, S. and Rheinberg, F. (2008). Flow, performance and moderators of challenge-skill balance. \emph{Motivation and Emotion}, 32(3).
\bibitem[Fahlman et~al.(2009)]{Fahlman2009} Fahlman, S.~A., Mercer, K.~B., Gaskovski, P., Eastwood, A.~E. and Eastwood, J.~D. (2009). Does a lack of life meaning cause boredom? Results from psychometric, longitudinal, and experimental analyses. \emph{Journal of Social and Clinical Psychology}, 28(3).
\bibitem[Fahlman et~al.(2013)]{Fahlman2013} Fahlman, S.~A., Mercer-Lynn, K.~B., Flora, D.~B. and Eastwood, J.~D. (2013). Development and validation of the Multidimensional State Boredom Scale. \emph{Assessment}, 20(1), 68--85.
\bibitem[Farmer and Sundberg(1986)]{FarmerSundberg1986} Farmer, R. and Sundberg, N.~D. (1986). Boredom proneness: The development and correlates of a new scale. \emph{Journal of Personality Assessment}, 50(1), 4--17.
\bibitem[Gerritsen et~al.(2014)]{Gerritsen2014} Gerritsen, C.~J., Toplak, M.~E., Sciaraffa, J. and Eastwood, J. (2014). I can't get no satisfaction: Potential causes of boredom. \emph{Consciousness and Cognition}, 27.
\bibitem[Goetz et~al.(2014)]{Goetz2014} Goetz, T., Frenzel, A.~C., Hall, N.~C., Nett, U.~E., Pekrun, R. and Lipnevich, A.~A. (2014). Types of boredom: An experience sampling approach. \emph{Motivation and Emotion}, 38(3).
\bibitem[Goldberg et~al.(2011)]{Goldberg2011} Goldberg, Y.~K., Eastwood, J.~D., LaGuardia, J. and Danckert, J. (2011). Boredom: An emotional experience distinct from apathy, anhedonia, or depression. \emph{Journal of Social and Clinical Psychology}, 30(6).
\bibitem[Havermans et~al.(2015)]{Havermans2015} Havermans, R.~C., Vancleef, L., Kalamatianos, A. and Nederkoorn, C. (2015). Eating and inflicting pain out of boredom. \emph{Appetite}, 85.
\bibitem[Hunter and Eastwood(2018)]{HunterEastwood2018} Hunter, A. and Eastwood, J.~D. (2018). Does state boredom cause failures of attention? Examining the relations between trait boredom, state boredom, and sustained attention. \emph{Experimental Brain Research}, 236(9).
\bibitem[Isacescu et~al.(2017)]{Isacescu2017} Isacescu, J., Struk, A.~A. and Danckert, J. (2017). Cognitive and affective predictors of boredom proneness. \emph{Cognition and Emotion}, 31(8).
\bibitem[Keller et~al.(2011)]{Keller2011} Keller, J., Bless, H., Blomann, F. and Kleinb{\"o}hl, D. (2011). Physiological aspects of flow experiences: Skills-demand-compatibility effects on heart rate variability and salivary cortisol. \emph{Journal of Experimental Social Psychology}, 47(4).
\bibitem[Larson and Richards(1991)]{LarsonRichards1991} Larson, R.~W. and Richards, M.~H. (1991). Boredom in the middle school years: Blaming schools versus blaming students. \emph{American Journal of Education}, 99(4).
\bibitem[Malkovsky et~al.(2012)]{Malkovsky2012} Malkovsky, E., Merrifield, C., Goldberg, Y. and Danckert, J. (2012). Exploring the relationship between boredom and sustained attention. \emph{Experimental Brain Research}, 221(1).
\bibitem[Mann and Cadman(2014)]{MannCadman2014} Mann, S. and Cadman, R. (2014). Does being bored make us more creative? \emph{Creativity Research Journal}, 26(2).
\bibitem[Mercer-Lynn et~al.(2014)]{MercerLynn2014} Mercer-Lynn, K.~B., Bar, R.~J. and Eastwood, J.~D. (2014). Causes of boredom: The person, the situation, or both? \emph{Personality and Individual Differences}, 56.
\bibitem[Merrifield and Danckert(2014)]{MerrifieldDanckert2014} Merrifield, C. and Danckert, J. (2014). Characterizing the psychophysiological signature of boredom. \emph{Experimental Brain Research}, 232(2).
\bibitem[Moneta and Csikszentmihalyi(1996)]{MonetaCsikszentmihalyi1996} Moneta, G.~B. and Csikszentmihalyi, M. (1996). The effect of perceived challenges and skills on the quality of subjective experience. \emph{Journal of Personality}, 64(2).
\bibitem[Mugon et~al.(2018)]{Mugon2018} Mugon, J., Struk, A. and Danckert, J. (2018). A failure to launch: Regulatory modes and boredom proneness. \emph{Frontiers in Psychology}, 9, 1126.
\bibitem[Nederkoorn et~al.(2016)]{Nederkoorn2016} Nederkoorn, C., Vancleef, L., Wilkenh{\"o}ner, A., Claes, L. and Havermans, R.~C. (2016). Self-inflicted pain out of boredom. \emph{Psychiatry Research}, 237.
\bibitem[O'Hanlon(1981)]{OHanlon1981} O'Hanlon, J.~F. (1981). Boredom: Practical consequences and a theory. \emph{Acta Psychologica}, 49(1).
\bibitem[Pekrun et~al.(2010)]{Pekrun2010} Pekrun, R., Goetz, T., Daniels, L.~M., Stupnisky, R.~H. and Perry, R.~P. (2010). Boredom in achievement settings: Exploring control--value antecedents and performance outcomes of a neglected emotion. \emph{Journal of Educational Psychology}, 102(3), 531--549.
\bibitem[Raffaelli et~al.(2018)]{Raffaelli2018} Raffaelli, Q., Mills, C. and Christoff, K. (2018). The knowns and unknowns of boredom: A review of the literature. \emph{Experimental Brain Research}, 236(9).
\bibitem[Sansone et~al.(1992)]{Sansone1992} Sansone, C., Weir, C., Harpster, L. and Morgan, C. (1992). Once a boring task always a boring task? Interest as a self-regulatory mechanism. \emph{Journal of Personality and Social Psychology}, 63(3).
\bibitem[Sommers and Vodanovich(2000)]{SommersVodanovich2000} Sommers, J. and Vodanovich, S.~J. (2000). Boredom proneness: Its relationship to psychological- and physical-health symptoms. \emph{Journal of Clinical Psychology}, 56(1).
\bibitem[Struk et~al.(2016)]{Struk2016} Struk, A.~A., Scholer, A.~A. and Danckert, J. (2016). A self-regulatory approach to understanding boredom proneness. \emph{Cognition and Emotion}, 30(8).
\bibitem[Struk et~al.(2017)]{Struk2017} Struk, A.~A., Carriere, J.~S.~A., Cheyne, J.~A. and Danckert, J. (2017). A short Boredom Proneness Scale: Development and psychometric properties. \emph{Assessment}, 24(3), 346--359.
\bibitem[Tam et~al.(2021)]{Tam2021} Tam, K.~Y.~Y., van Tilburg, W.~A.~P., Chan, C.~S., Igou, E.~R. and Lau, H. (2021). Attention drifting in and out: The Boredom Feedback Model. \emph{Personality and Social Psychology Review}, 25(3).
\bibitem[Tze et~al.(2016)]{Tze2016} Tze, V.~M.~C., Daniels, L.~M. and Klassen, R.~M. (2016). Evaluating the relationship between boredom and academic outcomes: A meta-analysis. \emph{Educational Psychology Review}, 28(1).
\bibitem[van Tilburg and Igou(2011)]{vanTilburgIgou2011} van Tilburg, W.~A.~P. and Igou, E.~R. (2011). On boredom and social identity: A pragmatic meaning-regulation approach. \emph{Personality and Social Psychology Bulletin}, 37(12).
\bibitem[van Tilburg and Igou(2012)]{vanTilburgIgou2012} van Tilburg, W.~A.~P. and Igou, E.~R. (2012). On boredom: Lack of challenge and meaning as distinct boredom experiences. \emph{Motivation and Emotion}, 36(2), 181--194.
\bibitem[van Tilburg and Igou(2016)]{vanTilburgIgou2016} van Tilburg, W.~A.~P. and Igou, E.~R. (2016). Going to political extremes in response to boredom. \emph{European Journal of Social Psychology}, 46(6).
\bibitem[van Tilburg and Igou(2017)]{vanTilburgIgou2017} van Tilburg, W.~A.~P. and Igou, E.~R. (2017). Boredom begs to differ: Differentiation from other negative emotions. \emph{Emotion}, 17(2).
\bibitem[van Tilburg et~al.(2013)]{vanTilburgIgouSedikides2013} van Tilburg, W.~A.~P., Igou, E.~R. and Sedikides, C. (2013). In search of meaningfulness: Nostalgia as an antidote to boredom. \emph{Emotion}, 13(3).
\bibitem[Vodanovich(2003)]{Vodanovich2003} Vodanovich, S.~J. (2003). Psychometric measures of boredom: A review of the literature. \emph{The Journal of Psychology}, 137(6).
\bibitem[Vodanovich and Kass(1990)]{VodanovichKass1990} Vodanovich, S.~J. and Kass, S.~J. (1990). A factor analytic study of the Boredom Proneness Scale. \emph{Journal of Personality Assessment}, 55(1--2).
\bibitem[Vodanovich and Watt(2016)]{VodanovichWatt2016} Vodanovich, S.~J. and Watt, J.~D. (2016). Self-report measures of boredom: An updated review of the literature. \emph{The Journal of Psychology}, 150(2).
\bibitem[Westgate and Wilson(2018)]{WestgateWilson2018} Westgate, E.~C. and Wilson, T.~D. (2018). Boring thoughts and bored minds: The MAC model of boredom and cognitive engagement. \emph{Psychological Review}, 125(5), 689--713.
\bibitem[Wilson et~al.(2014)]{WilsonEtAl2014} Wilson, T.~D., Reinhard, D.~A., Westgate, E.~C., Gilbert, D.~T., Ellerbeck, N., Hahn, C., Brown, C.~L. and Shaked, A. (2014). Just think: The challenges of the disengaged mind. \emph{Science}, 345(6192), 75--77.
\bibitem[Wolff et~al.(2020)]{Wolff2020} Wolff, W., Martarelli, C.~S., Sch{\"u}ler, J. and Bieleke, M. (2020). High boredom proneness and low trait self-control impair adherence to social distancing guidelines during the COVID-19 pandemic. \emph{International Journal of Environmental Research and Public Health}, 17(15), 5420.
\bibitem[Zuckerman(1979)]{Zuckerman1979} Zuckerman, M. (1979). \emph{Sensation Seeking: Beyond the Optimal Level of Arousal}. Lawrence Erlbaum.
\end{thebibliography}

\end{document}
